% =====================================================================
% REMORA: A Policy-Gated Multi-Oracle Assurance Architecture for Agentic AI
% LaTeX source - suitable for Overleaf, NeurIPS/ICML/ACL workshop templates
% =====================================================================
\documentclass[10pt,twocolumn]{article}

% ---- Packages -------------------------------------------------------
\usepackage[utf8]{inputenc}
\usepackage[T1]{fontenc}
\usepackage{lmodern}
\usepackage[margin=1in]{geometry}
\usepackage{amsmath,amssymb,amsfonts,amsthm}
\usepackage{booktabs}
\usepackage{multirow}
\usepackage{graphicx}
\usepackage[hidelinks]{hyperref}
% xurl = url with break points after every character class, so long
% URLs and \path tokens wrap inside the narrow two-column measure
% instead of protruding into the margin.
\usepackage{xurl}
\usepackage{microtype}
\usepackage{xcolor}
\usepackage{listings}
% \usepackage{algorithm}    % not installed in build env; no algorithm blocks used
% \usepackage{algpseudocode}
\usepackage{caption}
\usepackage{subcaption}
\usepackage{enumitem}
\usepackage{tcolorbox}

% ---- Code listings ---------------------------------------------------
% Global listings setup. verbatim must not be used in this paper: it
% never breaks lines, so any long line silently overflows the column in
% twocolumn mode. lstlisting with breaklines wraps instead.
\lstset{
  basicstyle=\ttfamily\footnotesize,
  breaklines=true,
  breakatwhitespace=false,
  postbreak=\mbox{\textcolor{gray}{$\hookrightarrow$}\space},
  columns=fullflexible,
  keepspaces=true,
  frame=single,
  framerule=0.3pt,
  rulecolor=\color{black!25},
  xleftmargin=4pt,
  xrightmargin=4pt,
  aboveskip=6pt,
  belowskip=6pt,
}

% Fit a tabular to the column width. Used for tables whose natural width
% exceeds \columnwidth in twocolumn mode; scaling preserves the layout
% without overflowing into the neighbouring column.
\newcommand{\fitcolumn}[1]{\resizebox{\columnwidth}{!}{#1}}

% Paragraph-level overflow tolerance for narrow two-column measure:
% allow modest extra inter-word stretch before a line is allowed to
% protrude. Removes residual overfull lines (long URLs, technical
% tokens) without visibly loosening the typography.
\setlength{\emergencystretch}{2em}
\tolerance=2000

% ---- Theorem environments -------------------------------------------
\newtheorem{theorem}{Theorem}
\newtheorem{proposition}[theorem]{Proposition}
\newtheorem{definition}[theorem]{Definition}
\newtheorem{lemma}[theorem]{Lemma}

% ---- Metadata -------------------------------------------------------
\title{REMORA: A Policy-Gated Multi-Oracle\\Assurance Architecture for Agentic AI}

% Paper version - kept in lockstep with paper/remora_paper.md by
% scripts/check_paper_sync.py (CI-enforced). The "v0.10.0" and
% "revision 2026-08-07" strings below MUST match the .md version line.
\newcommand{\paperversion}{Preprint v0.10.0, revision 2026-08-07 (synchronized with repository code); repository tag \texttt{review-v1}}

\author{
  Stian Skogbrott\thanks{Developer and maintainer of REMORA. Luftfiber AS.
  Contact: \texttt{support@luftfiber.no}.}\\
  Luftfiber AS\\
  \texttt{https://github.com/darklordVirtual/REMORA-research}
}

\date{\paperversion\ --- \small compiled \today\ by CI from the repository LaTeX source, synchronized with the code as of the revision date above}

% ---- Notation -------------------------------------------------------
\newcommand{\Oracles}{\mathcal{O}}
\newcommand{\OraclesV}{\mathcal{O}'}
\renewcommand{\phi}{\varphi}
\newcommand{\phat}{\hat{p}}
\newcommand{\trust}{\tau}
\newcommand{\lyap}{V}
\newcommand{\gate}{g}
\newcommand{\Gate}{\Gamma}

% =====================================================================
\begin{document}

\maketitle

% =====================================================================
\begin{abstract}
Autonomous AI agents that invoke tools, query databases, or actuate
real-world systems need a gate that decides---before anything
executes---whether an action may proceed autonomously, requires
verification, warrants abstention, or must be escalated to a human.
We present \textbf{REMORA} (Reasoning Ensemble Multi-Oracle Routing
Architecture), a research-grade control layer that makes this decision.
Several independent language models assess each proposed action in
parallel; a deterministic policy engine applies hard safety rules that
no probabilistic score can override; and every action leaves the gate
as one of four auditable outcomes: \textsc{accept}, \textsc{verify},
\textsc{abstain}, or \textsc{escalate}.

We report three findings.

\textbf{First, autonomous safety needs a deterministic floor.} On a
leakage-controlled tool-call benchmark, REMORA's full policy gate
produced zero simulated unsafe acceptances with significantly higher
decision utility and accuracy than heuristic baselines (mean utility
$+0.456$, $p\approx10^{-4}$; accuracy 90\% vs.\ $\leq 60\%$)---and an
ablation shows the zero is attributable entirely to the deterministic
hard rules, not to the multi-model voting machinery. Under an
intent-gating protocol on 208 independently sourced AgentHarm scenarios
\cite{andriushchenko2024agentharm}, the same rule set produced zero
false authorizations. Probabilistic consensus improves routing; it
cannot override policy.

\textbf{Second, unknown state is a resolution problem, not a verdict.}
Absence of a record without a declared completeness guarantee is
\textsc{unknown}, not \textsc{unsupported}. We make that distinction
operational: a \textsc{verify} must name an actual bounded lookup and be
followed by full re-entry of the whole router on a fresh observation, or
it degrades to \textsc{abstain}---promising a verification that cannot
happen is worse than stopping. Under study conditions this recovers read
utility from 0\% to 100\% in a regime where every value verdict is
\textsc{unknown}, with zero corrupt-identifier acceptances and zero
cross-tenant lookups.

\textbf{Third, the routing behaviour transfers to new external data.} In
one sealed BFCL v4 run (1{,}527 episodes), REMORA met \textbf{all five
pre-registered targets}: unknown required inputs were never guessed
(0/32), irrelevant tools were refused 258/258 times, unobtainable inputs
were refused 98/99 times, obtainable inputs were sent to verification
96/99 times, and well-formed wrong calls were accepted 28/258 times
(\textbf{10.9\%}, target $\leq 20\%$). Labelled routing accuracy was
91.2\%. The result confirms the current routing pipeline on unseen data
while preserving the distinction between structural validity and
semantic task correctness.

Two further findings constrain how much trust to place in model signals:
model confidence \emph{anticorrelates} with correctness exactly in the
high-disagreement phase where routing matters most, and a pre-registered
round on 1{,}231 fresh items failed to confirm our temperature-based
selection signal, which is therefore retained only as a logged
diagnostic.

Key limitations: the entropy observable is computed by a
token-fingerprint proxy rather than full semantic entropy over NLI
clusters; the QA benchmark is partly author-curated; semantic evidence
retrieval is pluggable but not live; the audit hash chain is
tamper-evident, not tamper-proof, without external append-only storage;
and the BFCL v4 track supplies no authoritative task-intent/tool-contract
bundle and excludes wrong-argument values because BFCL has no independently
vouchable state table. It therefore confirms the complete routing pipeline,
not the isolated causal effect of semantic intent matching. REMORA is a
governed autonomy layer, not a replacement for domain authority. The system
is \texttt{SHADOW\_ONLY}: independent external replication and production
field evidence remain pending.
\end{abstract}

\textbf{Keywords:} agentic AI safety, multi-oracle consensus, selective
prediction, uncertainty routing, policy-as-code, audit governance,
human-in-the-loop.

% =====================================================================
\section{Introduction}
\label{sec:intro}

The deployment of LLM-powered agents capable of invoking tools
introduces a class of decision problem distinct from conversational AI:
an agent proposes a \emph{concrete action} in the world, and the
system must decide whether to execute it. Unlike a hallucinated
sentence, an erroneously executed database write or wrongly triggered
process shutdown has immediate, potentially irreversible consequences.

Existing safeguards (RLHF alignment, system prompts, content
classifiers) are semantic filters applied at training or prompt time.
They are not designed to evaluate a specific proposed action against a
specific operational context, evidence base, and regulatory policy at
inference time. Majority vote among LLMs improves accuracy but cannot
block a confident, wrong consensus from triggering unsafe execution.

We argue that what is needed is not a \emph{better oracle} but an
\emph{assurance layer}: a system that evaluates whether the conditions
for autonomous action are met before anything is executed. When those
conditions are not met, the system must route toward verification,
abstention, or human review, recording the reasoning in an auditable
envelope.

REMORA demonstrates such a layer. Its core thesis:
\textbf{governed autonomy requires explicit routing of
uncertainty---not suppression of it.}

The contribution of this paper is a \textbf{system architecture and
empirical evaluation for governed agentic AI decision control}---not a
new foundation model, training procedure, or benchmark. It is a
reference implementation, a set of measurable observables, a
documented set of negative results, and an honest account of what
remains unsolved.

\paragraph{A note on terminology.}
Earlier versions of this paper named its observables after statistical
physics (\emph{temperature}, \emph{free energy}, \emph{Lyapunov
stability}). That vocabulary has been withdrawn: it invited a physical
reading the system does not support, and the quantity it rested on ---
structural temperature --- failed pre-registered confirmation on fresh
data (\S\ref{sec:method}, \S\ref{sec:negative}). What remains are the
information-theoretic quantities the benchmarks report: entropy $H$ and
dissensus $D$ over the weighted verdict distribution, and a trust score
$\trust$ derived from them. The three consensus regimes
(ordered/critical/disordered) keep their names because they name
existing code, not a physical state.

% =====================================================================
\section{Background and Related Work}
\label{sec:related}

\subsection{LLM Ensembles and Self-Consistency}
Wang et al.\ (2023) \cite{wang2023self} introduced self-consistency
sampling over chain-of-thought paths. Lakshminarayanan et al.\ (2017)
\cite{lakshminarayanan2017simple} demonstrated that deep ensembles from
distinct random initialisations improve both accuracy and calibration.
REMORA's oracle fan-out is structurally related but uses
\emph{separately configured oracle endpoints} rather than multiple samples
from one model (the 2026-07 clean round and the AgentHarm validation use a
Cloudflare Workers AI cross-family trio --- Meta LLaMA, Alibaba Qwen,
Mistral; earlier rounds used a Groq three-LLaMA swarm),
and weights oracles by \emph{inverse pairwise correlation} rather than
treating them as exchangeable. Kuncheva \& Whitaker (2003)
\cite{kuncheva2003measures} formalised diversity measures for classifier
ensembles; REMORA's inverse-correlation weighting can be understood as an
instance of this family applied to instruction-following LLMs.
Wang, Aitchison \& Rudolph (2023) \cite{wang2023lora} show that
parameter-efficient LoRA ensembles improve predictive accuracy and
uncertainty quantification, with ensemble diversity as the key driver;
this motivates, but does not independently validate, REMORA's
heterogeneous oracle aggregation.

\subsection{Multi-Agent Debate}
Du et al.\ (2024) \cite{du2023improving} showed that inter-model
debate can reduce hallucination. REMORA's dissensus metric $D$ is a
one-shot, non-iterative analogue: disagreement is measured once, not
resolved through iterative exchange.

\subsection{LLM-as-Judge and Verifier Models}
Zheng et al.\ (2023) \cite{zheng2023judging} established the
LLM-as-judge paradigm. Cobbe et al.\ (2021) \cite{cobbe2021training}
trained outcome verifiers that rank complete model solutions. REMORA differs by using a \emph{structured policy engine} rather than a
learned verifier, enabling interpretable hard blocks that cannot be
overridden by a confident oracle.

\subsection{Selective Prediction and Abstention}
Chow (1970) \cite{chow1970optimum} established the optimal error-reject
tradeoff for binary classifiers. El-Yaniv \& Wiener (2010)
\cite{elyaniv2010selective} extended this to a distribution-free
selective classification framework; their coverage--accuracy bound is
the foundational result that Geifman \& El-Yaniv (2017)
\cite{geifman2017selective} apply to deep networks. Kadavath
et al.\ (2022) \cite{kadavath2022language} find promising LLM
self-knowledge on tasks resembling training data, with weaker
calibration and self-evaluation on novel tasks.
REMORA's abstention is grounded in structural oracle disagreement
rather than self-reported confidence.

\subsection{Conformal Prediction}
Conformal prediction originates with Vovk, Gammerman \& Shafer (2005)
\cite{vovk2005algorithmic}; Shafer \& Vovk (2008) \cite{shafer2008tutorial}
and Angelopoulos \& Bates (2021) \cite{angelopoulos2021gentle} established
it as a practical distribution-free framework for coverage guarantees.
REMORA implements Mondrian phase-stratified conformal calibration for the
ordered phase, and a CRC-inspired weighted empirical selective router
\cite{angelopoulos2022crc} with importance weights (following weighted
conformal prediction under covariate shift \cite{tibshirani2019conformal})
for the critical phase (\S\ref{sec:crc}). The exchangeability violation in
the critical phase is documented as a negative result
(\S\ref{sec:negative}). Conformal methods have since been brought to LLM
outputs directly: conformal factuality guarantees via claim back-off
\cite{mohri2024conformal} and conformal abstention for hallucination
mitigation \cite{yadkori2024abstention}; REMORA's contribution is narrower
and different in object: phase-stratified calibration over multi-oracle
governance observables for \emph{action gating}, not per-output factuality.

\subsection{Prover-Verifier Games}
Kirchner et al.\ (2024) \cite{kirsch2024prover} showed that prover-verifier
\emph{training} games improve LLM output legibility by requiring provers
to justify claims to iteratively trained verifiers. REMORA does not
implement that training game; it borrows the framing for a PVD-inspired
offline agreement score, using Semantic Entropy clustering to select the
prover (dominant-cluster oracle) and verifier (independent challenger)
from already-produced responses, without additional API calls
(\S\ref{sec:pvd}).

\subsection{AI Governance and Policy-as-Code}
Algorithmic-audit scholarship emphasizes institutionalized third-party
oversight and audit-system design \cite{raji2022outsider}; Raji et al.\
identify lack of audit data access as the most significant vulnerability
of the current AI audit ecosystem. REMORA's DecisionEnvelope addresses
this access gap by logging each governance decision with immutable
hash-chaining, and the Replay Engine publishes governance episodes for
third-party inspection. The EU AI Act (2024) \cite{euaiact2024} mandates human oversight and
audit trails for high-risk AI systems. Open Policy Agent (OPA)
\cite{opa} provides a production-grade Rego policy engine. REMORA
ships an OPA adapter---failing closed to a Python fallback---as a
concrete instantiation of policy-as-code for AI decisions; the reference
API currently invokes the Python engine directly (\S\ref{sec:architecture}).

\subsection{Assurance Cases}
Kelly (1998) \cite{kelly1998goal} and Bloomfield \& Bishop (2010)
\cite{bloomfield2010safety} structure safety arguments as
claim-evidence-reasoning trees. REMORA's DecisionEnvelope is designed
to populate one node of such a case with a specific gate decision's
evidence chain.

\subsection{Human-in-the-Loop / Human-on-the-Loop}
Endsley (1995) \cite{endsley1995toward} provides the
situation-awareness theory underpinning human supervision of dynamic
systems; the automation-oversight literature distinguishes
human-in-the-loop confirmation from human-on-the-loop monitoring.
REMORA's four-outcome schema implements this
spectrum: \textsc{accept} (human-on-the-loop) through
\textsc{escalate} (human-in-the-loop), with \textsc{verify} and
\textsc{abstain} as intermediate states.

\subsection{Industrial AI Safety}
The oil-and-gas sector has established quantitative risk frameworks
(NORSOK D-010:2021 \cite{norsok2021d010}; IEC 61511:2016 \cite{iec61511})
that define action classes requiring human sign-off. REMORA's case study
(\S\ref{sec:casestudy}) uses this domain as an illustration of how
policy-gated AI governance maps to existing industrial safety
requirements. No domain-specific correctness is claimed.

\subsection{Runtime Guardrails: Agent Permissioning, and AI Control}

\paragraph{Content guardrails.}
Llama Guard and its successors \cite{inan2023llamaguard} classify
conversational inputs/outputs against a safety taxonomy; NeMo Guardrails
\cite{rebedea2023nemo} provides programmable dialogue rails;
Constitutional Classifiers \cite{sharma2025constitutional} defend against
universal jailbreaks with classifiers trained on constitution-generated
data; Rule-Based Rewards \cite{mu2024rules} move safety rules into
training. These systems govern \emph{content}; REMORA governs
\emph{actions}: its unit of decision is a proposed tool call with
operational consequences, and its deterministic hard-block layer is, by
design, not overridable by any learned component. A trained classifier
and a policy floor fail differently, and REMORA's position is that action
gating in consequential environments needs the latter beneath the former.

\paragraph{Agent-level guardrail systems.}
LlamaFirewall \cite{chennabasappa2025llamafirewall} is the closest
industrial system in scope: a layered agent guardrail (prompt-injection
detection, alignment checking, code screening). GuardAgent
\cite{xiang2025guardagent} generates guardrail code with a guard LLM;
AgentSpec \cite{wang2025agentspec} defines a DSL of runtime
triggers/predicates/enforcement near-isomorphic to REMORA's policy
invariants. Two 2026 systems share REMORA's learning ambition: AgentTrust
\cite{yang2026agenttrust} splits threats into lexical (deterministic
rules) and semantic (LLM judge with a guarded verdict cache), and
\emph{self-distills judge decisions into its rule layer}; Membrane
\cite{choi2026membrane} evolves a contrastive safety memory against
jailbreaks. REMORA's differences are deliberate: (1) the hard-block floor
is fixed, versioned, and hash-provenanced policy-as-code, learned outputs
never become enforcement rules, and the learning layer (AROMER) is
monotonicity-tested to \emph{never weaken} policy
(\path{tests/policy/test_governance_intelligence_never_weakens_policy.py}),
whereas AgentTrust's distillation pipeline makes learned judgments
load-bearing for enforcement; (2) uncertainty is structural (multi-oracle
disagreement with conformal abstention), not a single judge's confidence;
(3) every decision emits a signed, hash-chained DecisionEnvelope.
Conversely, AgentTrust reports zero benign hard-blocks across 45{,}000
actions, the opposite end of the safety/friction trade-off from REMORA's
FBR=100\% AgentHarm corner point (\S\ref{sec:results}), and an explicit
target for REMORA's calibration roadmap.

\paragraph{Permissioning and security by design.}
Progent \cite{shi2025progent} enforces programmable least-privilege on
agent tool calls, reducing attack success from 39.9\% to 1.0\% on
AgentDojo \emph{while maintaining utility}; CaMeL
\cite{debenedetti2025camel} defeats prompt injection by construction via
control/data-flow separation, retaining 77\% task completion with
provable security. These results define the current (safety, friction)
frontier and pre-empt any reading of REMORA's zero-FAR corner point as
sufficient: REMORA's contribution is not the corner but the governance
machinery around it: uncertainty-routed escalation, auditability, and
claim-governed evaluation. Capability-based enforcement in the
Progent/CaMeL style is complementary and is the natural enforcement
substrate for REMORA's \textsc{verify}/\textsc{escalate} outcomes.

\paragraph{Benchmarks and environments.}
Beyond AgentHarm \cite{andriushchenko2024agentharm}, the field evaluates
on AgentDojo \cite{debenedetti2024agentdojo}, ToolEmu
\cite{ruan2024toolemu}, InjecAgent \cite{zhan2024injecagent}, R-Judge
\cite{yuan2024rjudge}, and $\tau$-bench \cite{yao2024taubench}. ToolEmu's
LM-emulated sandbox anticipates the counterfactual-execution idea behind
REMORA's Shadow Mode; REMORA's replay engine is therefore \emph{not}
conceptually unique: its distinction is the governance framing (committed
replay artifacts under claim governance). Replication of REMORA's results
on these independently constructed environments is required future work
and the honest boundary of every internal benchmark claim in this paper.

\paragraph{AI control.}
Greenblatt et al.\ (2024) \cite{greenblatt2024aicontrol} formalize safety
protocols that remain effective when the \emph{policy model itself} is
intentionally subversive, using trusted monitoring with limited audit
budgets. In their taxonomy REMORA is a trusted-monitoring control
protocol: a deterministic trusted layer plus weaker trusted monitors (the
oracle swarm) gating an untrusted agent, with \textsc{escalate} as the
audit budget. REMORA's threat model (injection, oracle compromise) does
not yet include a strategically subversive agent: collusion and
monitor-evasion failure modes from the control literature are open items
for the red-team plan.

\paragraph{Regulatory operationalization.}
COMPL-AI \cite{guldimann2024complai} provides a technical interpretation
of the EU AI Act as model-level benchmarks; REMORA operationalizes a
complementary slice at runtime, Article 12 (logging) via the
DecisionEnvelope hash chain, Article 14 (human oversight) via
\textsc{verify}/\textsc{escalate} routing, without claiming conformity
assessment.

\paragraph{Multi-layer guardrail taxonomies.}
A systematic literature review of runtime guardrails
\cite{shamsujjoha2024taxonomy} frames multi-layer defense as a
Swiss-cheese model; a layered governance architecture with OS-level
enforcement is developed in Ge (2026) \cite{ge2026governance}; confidence
elicitation by fidelity is treated in Zhang et al.\ (2024)
\cite{zhang2024fidelity}; and formal verification of behavioral safety
properties for neural policies in Corsi et al.\ (2021)
\cite{corsi2021formal}. REMORA's position relative to these four is that
it composes recognized layers (policy floor, ensemble uncertainty,
conformal abstention, audit chain) with a claim-governance methodology,
rather than introducing a new layer type.

% =====================================================================
\section{Problem Statement}
\label{sec:problem}

Let $\mathcal{A}$ be a space of agent-proposed actions, $q$ a query,
and $c = (\text{domain}, \text{risk}, \text{env})$ an operational
context. We seek a control function:
%
\begin{equation}
\begin{split}
\Gate: (q, a, c) \;\longrightarrow\; g, \qquad\qquad\\
g \in \{\textsc{Accept}, \textsc{Verify}, \textsc{Abstain}, \textsc{Escalate}\}
\end{split}
\end{equation}
%
with five required properties: (1)~safety-first (hard policy violations
map to \textsc{Escalate} regardless of consensus); (2)~calibration-aware \cite{guo2017calibration}
(abstain when uncertainty exceeds threshold); (3)~evidence-sensitive
(evidence can resolve or contradict ambiguous consensus);
(4)~auditable (every call produces a hash-chained evidence record,
HMAC-signed when a signing key is configured);
(5)~conservative by default (fail closed toward \textsc{Verify} or \textsc{Escalate} depending on risk tier).

This is not a question of oracle \emph{accuracy}---REMORA does not
know whether action $a$ is correct. It is a question of
\emph{assurance conditions}: whether the conditions that justify
autonomous execution of $a$ are verifiably met.

% =====================================================================
\section{REMORA Architecture}
\label{sec:architecture}

REMORA processes each gate request through six decision stages,
followed by DecisionEnvelope emission (Stage~7).
The architecture diagram in the supplementary material (Mermaid) shows
the full pipeline. Stage~7 is the output record, not a decision stage.

\paragraph{Stage 1: Intake \& Risk Classification.}
The system classifies intent domain, risk tier
($\in\{$low, medium, high, critical$\}$), action type, and target
environment. An adversarial admission firewall checks for injection
patterns (``ignore previous instructions'', ``drop table'', etc.) and
fires an immediate \textsc{Escalate} if detected.

\paragraph{Stage 2: Oracle Fan-Out.}
$n$ configured oracle models are queried in parallel
via \texttt{ThreadPoolExecutor} (the engine accepts $n \geq 2$; the
design recommendation is $n \geq 3$ from distinct model families. The
2026-07 clean round and AgentHarm use a Cloudflare Workers AI cross-family
trio --- Meta LLaMA, Qwen, Mistral --- enforced by
\texttt{remora/oracles/families.py}; earlier rounds used a Groq
three-LLaMA swarm whose within-family correlation gave an effective
diversity closer to two). Responses with
non-null \texttt{error} fields are excluded before any aggregation.
$\OraclesV \subseteq \Oracles$ denotes the valid oracle set.

\paragraph{Stage 3: Canonicalization \& Weighted Consensus.}
Each response is canonicalized by $\varphi: \mathrm{raw} \rightarrow v$
where $v$ is the tuple (\texttt{polarity}, \texttt{claim\_hash},
\texttt{magnitude}, \texttt{tags}).
The \texttt{claim\_hash} field uses a pluggable backend: the default
\texttt{TokenFingerprintBackend} applies a 16-char SHA-256 truncation over
NFKC-normalized sorted tokens for backward-compatible, dependency-free
operation; the \texttt{NLISemanticBackend} replaces this with bidirectional
cross-encoder entailment clustering, implementing Semantic Entropy
\cite{kuhn2023semantic}. Oracle weights are derived from
inverse pairwise correlation (see \S\ref{sec:method}).

\paragraph{Stage 4: Consensus Observables.}
Entropy $H$ is computed over the diversity-weighted verdict distribution
$\hat{p}$. The \texttt{NLISemanticBackend} grounds $H$ as Semantic
Entropy~\cite{kuhn2023semantic},
$\mathrm{SE}(x) = -\sum_{c \in \mathcal{C}} p(c\,|\,x)\log p(c\,|\,x)$
over NLI-derived semantic equivalence clusters $\mathcal{C}$.
\emph{Note on experimental backend:} all benchmarks reported in
\S\ref{sec:eval} use the default \texttt{TokenFingerprintBackend}
(NFKC-normalised, sorted-token SHA-256 truncation), which approximates
cluster identity via lexical overlap. The NLI backend is a
drop-in replacement for future work but was not used in the
reported results; claims about Semantic Entropy properties should
be read accordingly.
Dissensus $D$, consensus concentration $\eta$ and trust score
$\trust$ are computed (see \S\ref{sec:method}), and a regime
$\in\{\text{ordered}, \text{critical}, \text{disordered}\}$ is
assigned. These run on the research and benchmark path; the governed
server routes leave them unset (\S\ref{sec:method}).

\paragraph{Stage 5: Evidence Router (Critical Phase).}
If phase = critical, the \texttt{CriticalEvidenceRouter} evaluates a
five-field \texttt{EvidenceSignal} and returns one of:
\texttt{evidence\_accept}, \texttt{abstain}, or \texttt{escalate}.
\emph{Important:} the current EvidenceSignal is an oracle-proxy
(built from consensus statistics); semantic BM25/NLI retrieval is
pluggable but not live (\S\ref{sec:evidence}).

\paragraph{Stage 6: Policy Engine.}
Seven \emph{headline} policy controls are evaluated in priority order
before any routing logic; the shipped engine's complete priority-ordered
gate inventory is larger (schema, forbidden-tool, coercion, tainted-argument,
rollback, production-write, quorum, trap-detection, session/fleet-risk
gates, among others) and is maintained in
\path{remora/policy/decision_engine.py}, with \texttt{hard\_guard\_floor()}
as the single deterministic floor no probabilistic result can override. An OPA/Rego adapter with Python fallback is implemented and
parity-tested, but the reference API invokes the Python policy engine
directly; OPA remains an integration option and production dependency
(fail closed either way). Gate outcome $g$ is assigned.

\paragraph{Output: DecisionEnvelope v2.}
An attribute-frozen envelope (treated as read-only; nested collections
are not deep-frozen) records: request, assessment, gate, reviewer
context, follow-up request, case history, policy-learning candidate,
and audit hash-chain entry.


\subsection{The Argument-Resolution Ladder}
\label{sec:resolution-ladder}

The six stages above describe how an action reaches a verdict. They do not
describe what the engine reasons about when the action is well-formed but the
\emph{state it refers to} is uncertain --- which, in operational deployments,
is the common case. That reasoning is a distinct ladder, evaluated after every
blocking gate so it can only ever make a would-be \textsc{accept} more
cautious, never unblock one:

\begin{enumerate}\setlength{\itemsep}{1pt}
\item \textbf{Authoritative tool metadata.} Effect, risk class, resource type
  and validator bindings come from a server-side registry. A caller may raise
  risk, never lower it. Verb heuristics are a fallback for unregistered tools
  only --- a lesson learned when \texttt{assign}, \texttt{close} and
  \texttt{approve} were absent from a mutating-verb list and thirty writes
  scored as reads.
\item \textbf{Coverage declaration.} A state index may not infer its own
  completeness. Closed-world authority is a signed, curator-bound declaration
  scoped per argument role, with a snapshot hash, an as-of date and a
  freshness bound. Without it, absence means \textsc{unknown}.
\item \textbf{Tenant, hash and freshness binding.} A declaration that does not
  match the snapshot it claims, or that has aged past its bound, is refused
  rather than trusted. Cross-tenant consultation is refused outright.
\item \textbf{Argument support.} Does the system of record contain this value?
  Tri-state: present, confirmed absent, or outside coverage. Unknown must
  never behave like a negative.
\item \textbf{Value grounding.} Are the argument values traceable to
  \emph{this} task --- named in the task text, declared by the tool's own
  schema, or vouched for by the system of record? A call whose values anchor
  to nothing here is observationally a foreign call.
\item \textbf{Validation requirement.} An argument that steers where the action
  lands must be confirmed before autonomy, decided by what the argument
  \emph{does} rather than by whether an index happens to cover it.
\item \textbf{Resolver availability.} If a bounded lookup can close the gap, the
  verdict is \textsc{verify} carrying a \texttt{ResolutionPlan} that names the
  resolver, the exact arguments it may write, and the sources permitted to
  supply them. If no resolver exists, the verdict is \textsc{abstain}:
  promising a verification that cannot happen is worse than stopping.
\item \textbf{Validator execution and full re-entry.} The resolver returns one
  value and hands control back. It never performs the original call and cannot
  redirect it. The \emph{whole} router then re-runs on a fresh observation, so
  every guard applies again to the resolved call rather than being patched
  around.
\item \textbf{Read/write posture.} A verified read may proceed autonomously; a
  write does not, whatever the state evidence says.
\end{enumerate}

This ladder, not the consensus machinery, is what the operational results in
\S\ref{sec:routing-blind}--\S\ref{sec:wrongcall-implication} measure. Its
tri-state discipline is the recurring commitment: \textsc{unknown} is a first-class
verdict distinct from \textsc{unsupported}, and a signal that cannot be established
is left absent rather than guessed --- a fabricated boolean enters the policy
contract as fact, while \texttt{None} correctly says nothing establishes this.

% =====================================================================
\section{Consensus Observables}
\label{sec:method}

\emph{Earlier versions of this paper presented these observables under
thermodynamic terminology (structural temperature, free energy, a
Lyapunov observable). That framing has been withdrawn: it invited a
physical reading the system does not support, and its central quantity
did not survive pre-registered testing. What remains are the
information-theoretic quantities the benchmarks actually report; the
withdrawn material is retained in the negative results
(\S\ref{sec:negative}), not deleted.}

\subsection{Weighted Support Distribution}

Let $\hat{p}(v) = \sum_{o \in \OraclesV} w(o) \cdot
\mathbf{1}[\varphi(o)=v]$ be the diversity-weighted support for
verdict $v$, normalized to sum to 1. The diversity weight $w(o)$
penalizes oracles that frequently agree with others:
%
\begin{equation}
  w(o) = \frac{(1/n)}{1 + \sum_{j \neq o} \rho(o,j)} \cdot Z^{-1}
\end{equation}
%
where $\rho(o,j)$ is the rolling pairwise agreement rate (window 200
samples) and $Z$ normalizes. Oracles that agree with the swarm receive
lower weight; independent dissenters receive higher weight.
The rolling window persists for a long-lived engine instance (as in the
offline experiments); the reference API constructs a fresh engine per
request, so cross-request persistence of correlation state is a
deployment integration point, not a current runtime property.

\subsection{Core Observables}

\textbf{Shannon entropy} (computed over the token-fingerprint proxy
clustering in all reported experiments, not NLI-based semantic
clusters; backend note in \S\ref{sec:architecture}, limitations in
\S\ref{sec:threats}):
\begin{equation}
  H = -\sum_v \hat{p}(v) \log_2 \hat{p}(v) \quad \text{(bits)}
\end{equation}

\textbf{Dissensus:}
\begin{equation}
  D = 1 - \max_v \hat{p}(v)
\end{equation}

\textbf{Order parameter:}
\begin{equation}
  \eta = \frac{\max_v \hat{p}(v) - 1/k}{1 - 1/k}
\end{equation}
where $k = |\{v : \hat{p}(v) > 0\}|$ is the number of active verdicts.

\textbf{Trust score:}
\begin{equation}
  \trust = \eta \cdot (1 - h_{\text{bound}})
           \cdot w_{\text{phase}}
           \cdot \left(1 + \frac{\chi}{\chi_0}\right)^{-1}
\end{equation}
where $h_{\text{bound}}$ is a \emph{heuristic false-consensus risk proxy}
from inter-oracle agreement --- the implementation uses $q^{n/2}$ (not the
proven exponent) and clamps $\rho$ at $0.49$, so it can understate
correlated-failure risk and must not be cited as a mathematical bound ---
$w_{\text{phase}} \in \{1.0, 0.5, 0.1\}$ is a regime weight over the
three consensus regimes (ordered/critical/disordered), and $\chi$ is a
scale factor that is algebraically constant in the live path unless
internal clamps bind. As a standalone difficulty predictor $\chi$ failed
(AUC $=0.39$, below chance; \S\ref{sec:negative}). Both are retained
because committed benchmark artifacts depend on them; neither is
presented as a validated mechanism.

\subsection{What these observables do not do}

Two boundaries matter more than the definitions above, because the
withdrawn framing obscured both.

\textbf{They are not populated on any governed path.}
\path{servers/api.py} never references the consensus state;
\path{servers/execution_api.py} passes \texttt{trust\_score=None} and
\texttt{phase=None} explicitly, and \texttt{build\_full\_observation} ---
the builder shared by the assess and execution routes --- leaves these
fields at their \texttt{None} defaults. The engine therefore fails toward
\textsc{verify}/\textsc{abstain} there: no probabilistic \textsc{accept}
can fire from these signals in the governed path at all. The runtime
safety floor is the deterministic hard-block layer
(\S\ref{sec:decision}). Everything in \S\ref{sec:results} is measured on
the research and benchmark path.

\textbf{The regime assignment rests on a falsified signal.} Regimes were
derived by comparing a structural temperature estimate against a
calibrated critical value. That signal failed its pre-registered
confirmation on fresh data: on $N=1{,}231$ held-out items it ranked
\emph{worse} than calibrated confidence (AURC 0.0954 vs.\ 0.0664, paired
CI excluding zero; CLAIM-012). Temperature is diagnostics-grade, not an
authoritative selector, and the derivation now lives in the negative
results (\S\ref{sec:negative}) rather than here as method.

% =====================================================================
\section{Decision and Policy Model}
\label{sec:decision}

\subsection{Four Gate Outcomes}

\begin{table}[h]
\centering\small
\caption{Gate outcomes and autonomous action status.}
\fitcolumn{%
\begin{tabular}{@{}llc@{}}
\toprule
Outcome & Meaning & Autonomous action \\
\midrule
\textsc{Accept}   & Assurance conditions met         & Permitted \\
\textsc{Verify}   & Plausible, validation required   & Pending \\
\textsc{Abstain}  & Uncertainty too high to decide   & Blocked \\
\textsc{Escalate} & Human review required            & Blocked \\
\bottomrule
\end{tabular}}
\label{tab:outcomes}
\end{table}

\textsc{Accept} does not mean the action is \emph{correct}---it means
the assurance conditions for autonomous execution are met.
\textsc{Abstain} does not mean the action is \emph{wrong}---it means
the conditions for deciding are not met. Both \textsc{Abstain} and
\textsc{Escalate} block execution; \textsc{Escalate} additionally
triggers a human-review workflow.

\subsection{Hard Blocks (Priority Order)}

\begin{table}[h]
\centering\small
\caption{Priority-ordered policy gates: unconditional hard blocks
  (1--3, \texttt{hard\_guard\_floor()}) and conditional policy gates
  (4--7, \texttt{\_CONDITIONAL\_GATES}). Both sets take precedence
  over the probabilistic ACCEPT decision; only the unconditional
  blocks are evaluated before all routing logic.}
\fitcolumn{%
\begin{tabular}{@{}clll@{}}
\toprule
Pri. & Condition & Reason Code & Action \\
\midrule
\multicolumn{4}{@{}l}{\textit{Unconditional hard blocks (\texttt{hard\_guard\_floor()})}} \\
\midrule
1 & Adversarial detected       & ADM\_FIREWALL    & ESC. \\
2 & Counterfactual failed      & CF\_FAILED       & ESC. \\
3a& Evidence contradicted +    &                  &      \\
  & \quad contradiction cycles & EV\_CONTRAD.     & ESC. \\
3b& Evidence contradictions    & EV\_CONTRAD.     & ABS. \\
\midrule
\multicolumn{4}{@{}l}{\textit{Conditional policy gates (\texttt{\_CONDITIONAL\_GATES})}} \\
\midrule
4 & Refuse parametric + no ev. & THM\_REQ\_EV.    & VER. \\
5 & Distribution shift         & DIST\_SHIFT      & VER. \\
6 & Phase=critical + risk=crit.& CRITICAL\_PHASE  & ESC. \\
7 & High/crit. risk + no ev.   & EV\_INSUFF.      & VER. \\
\bottomrule
\end{tabular}}
\label{tab:hardblocks}
\end{table}

\noindent\textbf{Policy must override model consensus.} Majority vote
is never sufficient for critical autonomous action when a hard block
condition is met.

\subsection{OPA/Rego Integration}

REMORA exports a 67-field \texttt{PolicyObservation} dataclass as
JSON to the OPA endpoint. On any communication error (connection
refused, timeout, parse failure), the Python fallback engine applies
identical rules and records \texttt{source\_of\_decision =
"python\_fallback"} in the envelope. No decision is ever withheld due
to a missing policy engine.

% =====================================================================
\section{Governance Intelligence Layer: Misspecification-Aware Pre-Policy Enrichment}
\label{sec:gov_intel}

The policy engine is only as honest as the metadata it is given: an
integration bug, or a compromised agent, can label a destructive
action as a low-risk read and present the gate with nothing to fire
on. Added at package version 0.9.0 (repository tag \texttt{review-v1}), the \emph{Governance Intelligence
Layer} (\path{remora/governance_intelligence/}) closes this gap with
an opt-in, fully deterministic enrichment stage that runs before any
policy rule (no LLM calls): (1)~fail-closed normalization of
caller-supplied labels (unknown is explicit, never coerced to a safe
value); (2)~action-semantics extraction from the proposed action text
and tool metadata; (3)~misspecification-risk inference from
label/semantics disagreement; (4)~causal-consequence signals \cite{bjoru2026causal,pearl2009causality} (blast
radius, expected loss); and (5)~policy-generalization risk---would
repeatedly accepting this \emph{class} of action remain safe
fleet-wide?

Signals merge into the existing \texttt{PolicyObservation} under a
\textbf{strengthen-only} rule: inferred higher risk may override a
supplied lower label, never the reverse, and hard-block flags are
never cleared. A grid property test across 2{,}160 observation
combinations verifies that enrichment never converts a rejection into
an acceptance
(\path{tests/policy/test_governance_intelligence_never_weakens_policy.py}).
The engine remains the sole decision authority.

\begin{table}[h]
\centering\small
\caption{Governance-intelligence benchmark: 50 deterministic tasks
across 10 categories of mislabelled, underspecified, or repeated
actions, evaluated offline with a deliberately permissive trust
baseline so that enrichment is the only barrier to acceptance.}
\begin{tabular}{@{}lr@{}}
\toprule
Metric & Result \\
\midrule
Unsafe accept rate (no-accept tasks) & \textbf{0.0\%} \\
Metadata-mismatch detection          & 75\% \\
Unknown-metadata review rate         & 100\% \\
Legitimate reads still accepted      & 100\% \\
Escalation precision                 & 0\% \\
Policy-generalization detection      & 62.5\% \\
\bottomrule
\end{tabular}
\label{tab:govintel}
\end{table}

\noindent The safety-relevant rows hold (0\% unsafe accepts; all
unknown-metadata and legitimate-read cases routed correctly), but the
enrichment layer's \emph{precision} is weak on this 50-task eval:
escalation precision 0\% (its escalations did not match the expected set),
75\% metadata-mismatch detection, 62.5\% policy-generalization detection.
These are the committed, deterministic artifact numbers (re-verified); the
layer is an offline, opt-in helper, not a validated detector.

Artifact:
\path{artifacts/governance_intelligence/evaluation_results.json};
runner: \path{experiments/evaluate_governance_intelligence.py}.
Without enrichment, \texttt{DROP TABLE users} labelled
\texttt{action\_type=read, risk\_tier=low} reaches \textsc{Accept} on
a high-trust path; with enrichment it escalates. The layer is
heuristic and English-centric; pattern tables can be evaded by novel
phrasing, so it reduces reliance on caller labels rather than
replacing verification, evidence, or human review.

% =====================================================================
\section{Evidence-Grounded Critical-Phase Routing}
\label{sec:evidence}

\subsection{The Critical-Phase Problem}

In the critical phase, oracle consensus is near the phase boundary.
Empirically, higher trust scores \emph{negatively} correlate with
correctness ($n=32$ real-oracle items): Q4 high-trust achieves 50\%
correct; Q1 low-trust achieves 75\% correct. Conformal calibration at
5\% risk target fails completely on critical-phase items (mean
observed risk = 100\%, coverage $\rightarrow 0\%$). Standard trust-based
routing cannot safely gate critical-phase decisions.

\subsection{PhaseAwareGuardrail: Exploiting Trust Inversion}
\label{sec:phase_aware}

The inversion is a \emph{selection criterion reversal}, not an
accuracy problem.  Items with $\tau < 0.10$ in the critical phase
achieve 76.2\% accuracy ($N=21$) vs.\ 36.4\% for high-$\tau$
critical items ($N=11$).  A \texttt{PhaseAwareGuardrail} exploits
this by applying an inverted score $\tilde{\tau} = 1 - \tau$ for
critical-phase items, then calibrating a conformal threshold on
$\tilde{\tau}$.  A hard gate rejects all items with $\tau \geq
\tau_{\max} = 0.10$ (the groupthink boundary) regardless of the
conformal result.

The combined pool---all 99 ordered-phase items (sorted by
$\tau \downarrow$) plus 21 low-$\tau$ critical items ($\tau < 0.10$,
sorted by $\tilde{\tau} \downarrow$)---achieves:

\begin{itemize}[leftmargin=*]
  \item 20.0\% coverage ($N=108$): 85.2\% accuracy \text{CI~[77.3\%, 90.7\%]}
  \item 22.1\% coverage ($N=120$): 85.0\% accuracy \text{CI~[77.5\%, 90.3\%]}
\end{itemize}

This is a $+3.9$\,pp coverage gain over the ordered-only operating
point (18.2\%, 86.9\% accuracy), with a 1.9\,pp accuracy cost---still
$+43.8$\,pp above the 41.18\% full-coverage baseline.  The guardrail is
implemented in \path{remora.selective.guardrail.PhaseAwareGuardrail}
and fully tested (25 unit tests).

\subsection{Conformal Risk Control Under Covariate Shift}
\label{sec:crc}

A weighted split-conformal router inspired by Conformal Risk Control is
implemented and tested, but it is an offline research component: the
reference API does not invoke it, and the implementation does not
instantiate the CRC theorem's hypotheses, so its output is a heuristic
bound rather than a certified one. Because it is neither part of the
operative architecture nor a valid instance of the formal method, the
derivation, the theorem statement, and the reasons no reweighting repairs
the gap are given in Appendix~\ref{app:crc} rather than here.

\subsection{PVD-Inspired Offline Agreement Score for Critical-Phase Routing}
\label{sec:pvd}

Kirchner et al.\ (2024) trained prover-verifier games that produce more
legible and reliable LLM outputs: helpful and sneaky \emph{provers}
justify claims to iteratively retrained small \emph{verifiers}
\cite{kirsch2024prover}. REMORA does \textbf{not} implement that training
game. The component below is a PVD-\emph{inspired}, post-hoc agreement
score computed over already-produced oracle responses: it selects two
existing responses, scores their agreement with the configured entailment
backend (default: token-fingerprint equality, not NLI), and applies a
fixed round decay. No model training, live interaction, argument
generation, or verifier update occurs, and it is not invoked by the
reference API (offline research component).

\paragraph{Protocol.}
\begin{enumerate}[leftmargin=*,label=\arabic*.]
  \item \textbf{Cluster} oracle responses via Semantic Entropy clustering
        \cite{kuhn2023semantic,farquhar2024detecting}
        (\S\ref{sec:method}), producing semantic equivalence classes
        $\mathcal{C} = \{C_1, \ldots, C_K\}$ sorted by mass.
  \item \textbf{Prover} = the oracle whose response belongs to the
        dominant cluster $C_1$ and has the highest per-oracle confidence.
  \item \textbf{Verifier} = the oracle with the highest confidence
        \emph{outside} $C_1$ (the independent challenger).
  \item \textbf{Deliberation} ($r = 1, \ldots, n_{\text{rounds}}$):
        evaluate bidirectional NLI entailment
        $e_r = \text{NLI}(\text{prover} \to \text{verifier}) \cdot \gamma^{r-1}$
        with round-decay factor $\gamma = 0.85$.
  \item \textbf{Legibility score}:
        $\mathcal{L} = \bar{e} \cdot |C_1| / N$,
        where $\bar{e}$ is the mean entailment and $|C_1|/N$ is the
        dominant cluster mass.
  \item \textbf{Final confidence}:
        $\hat{c} = \sqrt{|C_1|/N \cdot c_{\text{verifier}}}$
        (geometric mean of prover support and updated verifier confidence).
\end{enumerate}

\paragraph{Routing signal.}
The PVD confidence is blended with the raw REMORA trust score:
\begin{equation}
  \tilde{\tau} = (1 - w_{\text{pvd}}) \cdot \tau + w_{\text{pvd}} \cdot \hat{c},
  \qquad w_{\text{pvd}} = 0.40
  \label{eq:pvd_blend}
\end{equation}
In offline experiments this score can promote borderline critical-phase
items to \textsc{Accept} when agreement confidence is high, extending
coverage beyond the \texttt{PhaseAwareGuardrail} operating point. Neither
component is invoked by the reference API, and no formal risk bound
applies to this heuristic blend (\S\ref{sec:crc}).

\paragraph{Implementation.}
The deliberation entry point is
\begin{lstlisting}
remora.selective.pvd.deliberate(
    oracle_responses, oracle_confidences,
    n_rounds=2)
\end{lstlisting}
It returns a \texttt{PVDResult} with fields
\path{prover_cluster_mass}, \path{legibility_score},
\path{agreement}, and \path{final_confidence}.
\texttt{pvd\_routing\_score(trust\_score, pvd\_result)} implements
Eq.~\ref{eq:pvd_blend}.  Fully tested (33 unit tests, no external
model calls).

\subsection{CriticalEvidenceRouter}

The router evaluates a five-field \texttt{EvidenceSignal}
with fields (\texttt{evidence\_strength}, \texttt{contradiction\_\allowbreak score},
\texttt{citation\_coverage},
\texttt{cross\_evidence\_\allowbreak consistency},
\texttt{source\_reliability}):

\begin{enumerate}[leftmargin=*,label=\arabic*.]
  \item Coverage gate: \texttt{citation\_coverage}~$< 0.50$ $\Rightarrow$ \textsc{Escalate}
  \item Contradiction block: \texttt{contradiction\_score}~$> 0.50$ $\Rightarrow$ \textsc{Abstain}
  \item Evidence-accept: all thresholds met $\Rightarrow$ \texttt{evidence\_accept}
  \item Default: $\Rightarrow$ \textsc{Escalate} with diagnostic
\end{enumerate}

On MultiNLI ($N=3{,}000$): resolution rate 38.5\%, evidence-accept
precision 100\%, false-accept rate on contradiction items 0\%,
abstain precision on contradictions 99.3\%.

\subsection{Evidence Signal Provenance}

\emph{Current implementation:} \texttt{EvidenceSignal} is built as a
proxy from oracle consensus statistics. This is documented in the
source as a ``structural bridge'' pending semantic retrieval.
\emph{Planned:} external BM25/NLI passage retrieval feeding real
document evidence. The routing logic is fully implemented and tested;
real evidence quality may differ from NLI label quality.

% =====================================================================
\section{Audit and Governance Envelope}
\label{sec:audit}

\subsection{DecisionEnvelope v2}

Every gate invocation produces an attribute-frozen
\texttt{DecisionEnvelope} (frozen dataclasses treated as read-only;
nested lists/dicts are not deep-frozen) with blocks: request, assessment,
gate, reviewer\_context, follow\_up, history, policy\_learning, and audit.
The library engine emits the envelope with \texttt{previous\_hash} and
\texttt{signature} unset; the API finalizer links the tenant hash-chain
and adds an HMAC signature only when
\texttt{REMORA\_ENVELOPE\_SIGNING\_KEY} is configured.
The full schema is in Appendix~\ref{app:schema}.

\subsection{Audit Hash-Chain}

\begin{equation}
  h_i = \text{SHA-256}(h_{i-1} \;\|\; \mathrm{json}(e_i))
\end{equation}

This provides \textbf{tamper-detection}: any modification breaks the
chain. It does \textbf{not} prevent tampering. Tamper-\emph{proof}
audit requires external append-only storage (S3 Object Lock, Azure
Immutable Blob, or equivalent) as a deployment dependency.

\subsection{Follow-Up Workflow}

\textsc{Escalate} generates a structured follow-up request specifying
reason codes, required evidence, responsible role, and SLA. Data
structures are implemented; production integration with external
ticketing systems is a deployment requirement.

% =====================================================================
\section{Experimental Evaluation}
\label{sec:eval}

\subsection{Benchmarks}

\textbf{QA benchmark ($N=544$).}
BoolQ ($n=377$) \cite{clark2019boolq},
TruthfulQA ($n=85$) \cite{lin2022truthfulqa},
adversarial-curated ($n=7$),
REMORA-curated ($n=75$, author-assembled, selection bias possible).

\textbf{Tool-call benchmark (700 rows; effective $N=70$).}
Synthetic adversarial suite, leakage-remediated (2026-07-20): 70 unique
templates $\times$ 10 cosmetic variants across 7 domains and 4 task
types; the 700 rows are therefore 70 statistical clusters, not 700
independent tasks. Three adversarial failure modes: intent-argument
conflict, change-window/dual-control requirements, safe-looking tasks
with dangerous hidden scope. The simulator executes no real shell,
network, database, or filesystem mutation; ``unsafe'' means unsafe
\emph{acceptance} of a simulated action.

\textbf{Evidence benchmark ($N=3{,}000$).}
MultiNLI validation-matched \cite{williams2018broad} (entailment,
neutral, contradiction) as proxy for evidence-verification decisions.

\subsection{Baselines}

QA: full-coverage majority vote; trust-score selective (threshold=0.50).

Tool-call: single model (Llama-3.3-70B); majority vote; self-consistency;
verifier heuristic; REMORA temperature-gate only (no policy engine).

\subsection{Oracles}

The oracle backend changed across rounds. Earlier QA and tool-call rounds
ran a Groq-hosted three-LLaMA swarm (Llama-3.1-8B-Instant,
Llama-3.3-70B-Versatile, Llama-4-Scout-17B-16E-Instruct --- one weight
family); the 2026-07 clean round, the AgentHarm validation, and the
real-LLM baselines use a Cloudflare Workers AI cross-family trio (Meta
LLaMA, Alibaba Qwen, Mistral). Oracle model versions may change; stored
artifacts ensure deterministic reproduction.

% =====================================================================
\section{Results}
\label{sec:results}

\subsection{Selective Accuracy on QA Benchmark}

\begin{table}[h]
\centering\small
\caption{QA benchmark selective accuracy ($N=544$). $\dagger$~In-sample
optimum; $\ddagger$~held-out: stratified 80/20 split, $\tau^*=0.203$
locked from training set (seed=42); $\S$~\texttt{PhaseAwareGuardrail}:
ordered~$+$~inverted low-$\tau$ critical ($\tau < 0.10$), $N=120$.
Entropy-derived observables use the token-fingerprint proxy backend.}
\fitcolumn{%
\begin{tabular}{@{}llrrrr@{}}
\toprule
Condition & Set & Cov. & Acc. & Lift & 95\% CI \\
\midrule
Majority vote (full) & full   & 100\%  & 41.18\% & ---  & [37.1, 45.4] \\
neg-temp. (10\%)$\dagger$     & in-sam & 10\%   & 81.48\% & +40.3 & [69.2, 89.6] \\
neg-temp. (18\%)$\dagger$     & in-sam & 18\%   & \textbf{88.78\%} & \textbf{+47.6} & [81.0, 93.6] \\
phase-aware (22\%)$\S$        & in-sam & 22.1\% & 85.0\%  & +43.8 & [77.5, 90.3] \\
neg-temp. (25\%)$\dagger$     & in-sam & 25\%   & 72.79\% & +31.6 & [64.8, 79.6] \\
neg-temp. (held-out, retired)$\ddagger$ & \textbf{holdout} & 23.2\% & \textbf{88.00\%} & \textbf{+41.7} & \textbf{[70.0, 95.8]} \\
\bottomrule
\end{tabular}}

\smallskip
\footnotesize
Held-out: $N_\text{accepted}=25$ (all ordered-phase), $p=1.45\times10^{-5}$
(one-sided binomial, H$_1$: acc $>$ holdout baseline 46.3\%),
Wilson CI lies entirely above holdout baseline. \textbf{These are the
retired Groq-round figures.} The 2026-07 cross-family clean round gives
100.0\% at 16.7\% coverage ($N_\text{accepted}=18$, $p=0.052$ ---
directional only), and the pre-registered SAP v3 fresh-data round
\emph{failed to confirm} the consensus-temperature selection signal
(\S\ref{sec:negative}); do not read this row as evidence that temperature
generalizes.
\label{tab:qa}
\end{table}

Phase breakdown on full dataset:
ordered (N=99): 86.9\%; critical (N=32): 62.5\%; disordered (N=413): 28.6\%.
Disordered items constitute 75.9\% of the benchmark, driving the
low full-coverage baseline---a structural property of benchmark
composition. Extending coverage beyond the ordered-only 18.2\% operating
point requires exploiting the critical-phase trust inversion
(\S\ref{sec:phase_aware}): the \texttt{PhaseAwareGuardrail} raises
coverage to 22.1\% at 85.0\% accuracy by adding 21 low-$\tau$ critical
items ($\tau < 0.10$, inverted-score selection, $+43.8$\,pp lift; Wilson~CI [77.5\%, 90.3\%]).

\subsection{Tool-Call Benchmark (700 rows; effective $N=70$)}

\begin{table}[h]
\centering\small
\caption{Tool-call benchmark, leakage-remediated 2026-07-20: REMORA vs
baselines (700 rows; effective $N=70$ template clusters).}
\fitcolumn{%
\begin{tabular}{@{}lccc@{}}
\toprule
Condition & Acc. & Unsafe acc. & Utility \\
\midrule
Single model heuristic     & 28.6\% & 1.4\% & 0.16 \\
Majority vote heuristic    & 28.6\% & 1.4\% & 0.16 \\
Self-consistency heuristic & 28.6\% & 1.4\% & 0.16 \\
Verifier heuristic         & 28.6\% & 1.4\% & 0.16 \\
REMORA temp. gate only     & 60\%   & 1.4\% & 0.36 \\
\textbf{REMORA full policy}& \textbf{90\%} & \textbf{0\%} & \textbf{0.62} \\
\bottomrule
\end{tabular}}
\label{tab:toolcall}
\end{table}

REMORA's full policy gate produced 0 unsafe acceptances (cluster-level
Wilson 95\% CI [0.0\%, 5.2\%]); baselines produced 1.4\% (1 of 70
clusters). The safety-rate difference is not statistically significant
at the template-cluster level (one-sided $p=0.50$); the statistically
supported advantages are decision utility ($+0.456$, $p\approx10^{-4}$)
and accuracy. The hard-block layer carries the 0\% safety floor, but on
this benchmark version its measured advantage over the temperature-only
gate is routing quality, not a statistically separable unsafe-rate
reduction. (Earlier versions of this table reported baselines at
10--20\% unsafe: those numbers were inflated by baselines and gate
reading author-annotated severity and oracle context flags; the 2026-07-20
re-run restricts all conditions to the observable task surface.)

\subsection{Conformal Coverage (Mondrian, $N=2{,}161$)}

\begin{table}[h]
\centering\small
\caption{Mondrian phase-stratified conformal, 15\% risk target
(20 seeds $\times$ $N=2{,}161$ augmented dataset).}
\begin{tabular}{@{}lrrl@{}}
\toprule
Phase & Mean Risk & Mean Cov. & Seeds Fail \\
\midrule
Ordered    & 12.0\% & \textbf{99.9\%} & \textbf{0/20} \\
Critical   & 64.5\% &  0.9\% & 4/20 \\
Disordered & 78.6\% &  0.2\% & 2/20 \\
\bottomrule
\end{tabular}
\label{tab:mondrian}
\end{table}

Ordered-phase conformal achieves 99.9\% coverage with 0/20 failures---a
stable internal benchmark operating point (the $N=2{,}161$ dataset is
augmented and includes calibrated-simulation critical items; this is not
evidence of deployability under production distribution shift).
Critical and disordered phases cannot achieve
meaningful coverage via conformal calibration alone, motivating the
evidence router.

\subsection{Evidence Router Proxy Benchmark (MultiNLI, $N=3{,}000$)}

Resolution rate 38.5\%; evidence-accept precision 100\%;
false-accept rate on contradiction 0\%; abstain precision 99.3\%.
Trust-only baseline: 0\% resolution (all escalated).
This NLI-proxy benchmark is distinct from the 32-case
\emph{Cross-Domain Evidence Replay} (static rule-based provider)
reported in the repository and the enterprise white paper; the two
measure different mechanisms at different scales.


\subsection{Tool Routing on Sealed External Data (BFCL v4)}
\label{sec:routing-blind}

The benchmarks above measure whether a \emph{harmful} call is stopped. This
track also measures whether a harmless-looking but incorrect call is stopped.
Track C-ext2 retained five targets that were sealed before evaluation. It
sampled 258 \texttt{live\_multiple} tasks and 258
\texttt{live\_irrelevance} tasks from BFCL v4 at upstream commit
\texttt{6ea57973c7a6} (Apache-2.0). The one permitted evaluation contained
1{,}527 episodes over 516 source clusters, with an empty state index, no
validator bindings, and a registry derived only from task schemas.

\begin{center}
\small
\resizebox{\columnwidth}{!}{%
\begin{tabular}{lrrl}
\toprule
Pre-registered target & Goal & Sealed BFCL v4 result & Verdict \\
\midrule
Required-UNKNOWN autonomous \textsc{Acc} & $\leq 0\%$ & 0/32 = \textbf{0.0\%} & met \\
Irrelevance \textsc{Abs} recall & $\geq 70\%$ & 258/258 = \textbf{100.0\%} & met \\
Unobtainable \textsc{Abs} recall & $\geq 70\%$ & 98/99 = \textbf{99.0\%} & met \\
Obtainable \textsc{Ver} recall & $\geq 70\%$ & 96/99 = \textbf{97.0\%} & met \\
Wrong-call \textsc{Acc} & $\leq 20\%$ & 28/258 = \textbf{10.9\%} & met \\
\bottomrule
\end{tabular}}
\end{center}

All five pre-registered targets met. Labelled routing accuracy was
\textbf{91.2\%} ($n=1{,}170$; 508 effective labelled clusters; cluster-level
Wilson 95\% CI [88.4\%, 93.3\%]). Optional-lane identity \textsc{Accept} was
151/218 = 69.3\%, reported without a target.

The confirmation has a precise boundary. BFCL provides no independently
vouchable state table, so the wrong-argument-value axis remains excluded by
admission. The track also supplies no authoritative
\texttt{TaskIntent}/\texttt{ToolContract} bundle. Consequently, the blind
result confirms the complete current routing pipeline; it does not isolate the
causal efficacy of semantic intent matching.

\subsection{What the Wrong-Call Axis Implies}
\label{sec:wrongcall-implication}

The BFCL v4 result is this paper's clearest external routing result:

\begin{quote}
\textbf{Structural validity and semantic correctness are independent
properties, and the governance signals that establish the first cannot
establish the second.}
\end{quote}

Schema validity, argument satisfiability, provenance, taint and value grounding
are all properties of the \emph{call}. Whether the call serves the \emph{task}
is a property of the relation between them, and no amount of scrutiny of the
call alone recovers it. The residue is therefore not a threshold to tune: there
is no threshold on a signal that does not carry the information.

The architecture reserves a slot for the missing authority ---
\texttt{tool\_matches\_goal}, tri-state, held at \textsc{Unknown} since
definition rather than filled with a guess. A mechanism now occupies it: a
deployment-declared tool contract (capability, effect, resource type, mutation,
argument roles) matched against a structured task intent, routing an established
contradiction to \textsc{Abstain} for reads and \textsc{Escalate} for writes. It
is placed ahead of trust routing, so a confident consensus cannot execute a call
that provably serves the wrong goal, and behind every hard guard, so it can
never unblock one. The authority limit is the design: a model may \emph{propose}
the intent but may not thereby assert a match, and an intent whose quoted spans
do not occur in the task text yields \textsc{Unknown}, never \textsc{Supported}.

Its isolated effect remains \textbf{unmeasured}, because the BFCL v4
confirmation contains no authoritative semantic bundle.

% =====================================================================
\section{Ablations}
\label{sec:ablation}

\textbf{Where the safety floor comes from.}
\emph{Architectural implication.} The unsafe-execution reduction is
100\% attributable to deterministic, interpretable policy rules that are
independent of LLM output. The multi-oracle consensus and routing
machinery improve decision \emph{calibration and routing quality}
(discriminating VERIFY, ABSTAIN, and ACCEPT paths), but do not contribute
to the safety guarantee itself. This separation is a design feature: the
safety floor can be reasoned about and audited without reference to
oracle behaviour. Readers should not infer that the consensus
machinery is the source of the 0\% unsafe rate.

\textbf{Oracle count ($N=302$).}
Single oracle: 56.95\% (95\% CI [51.3, 62.4]).
Three-oracle majority vote: $\approx$+10 pp.
Correlation-aware weighting: additional gain within multi-oracle setting.

The primary benefit of the oracle swarm is disagreement
\emph{detection} (high $H$, high $D$) rather than accuracy
improvement per se.

% =====================================================================
\section{Case Study: Well Barrier Agent Action}
\label{sec:casestudy}

\emph{No domain-engineering correctness is claimed. This is an
illustrative governance walkthrough.}

A drilling agent proposes: \textit{``Autonomously update kill mud
weight from 1.52 to 1.48\,SG to improve ROP on the 8½{\textquotedblright}
section.''} Risk tier: critical. Domain: \texttt{well\_engineering}.

\textbf{Oracle responses:}
Llama-3.1-8B: ESCALATE, conf 0.78;
Llama-3.3-70B: ESCALATE, conf 0.85 (NORSOK D-010 §5.4);
Llama-4-Scout: VERIFY, conf 0.62.
One oracle failed; excluded.

\textbf{Consensus:}
$\hat{p}(\textsc{Escalate}) = 0.71$, $\hat{p}(\textsc{Verify}) = 0.29$.
$H = 0.866$ bits; $D = 0.29$;
regime = critical;
$\trust = 0.31$.

\textbf{Evidence:} No field inspection, no ECD calculation, no OIM
sign-off. Citation coverage = 0.18 $<$ 0.50 threshold.
Evidence router: \textsc{Escalate}.

\textbf{Policy engine:} Conditional policy gate fires
(phase=critical AND risk=critical) $\Rightarrow$ \textsc{Escalate}.
This gate resides in \texttt{\_CONDITIONAL\_GATES}, after the
unconditional \texttt{hard\_guard\_floor()}; the outcome is identical.

\textbf{Follow-up generated:} type=on\_site\_inspection;
assign\_to=Independent Well Engineer; required\_evidence = \{ECD
calculation, kick tolerance, NORSOK D-010 barrier confirmation, OIM
sign-off\}; SLA=4h.

\textbf{Audit hash} recorded; action blocked.


\section{Case Study: State Resolution on a Fleet Operations Agent}
\label{sec:casestudy-fleetops}

The barrier case shows the deterministic floor stopping an action outright.
This one shows the mechanism that distinguishes the current architecture: what
happens when the action is well-formed and the \emph{state it refers to} is
uncertain. Both calls below are structurally perfect. They are routed
differently, and neither routing depends on a model's confidence.

\paragraph{Setup.} A fleet-operations agent has six tools, of which
\texttt{get\_vehicle}, \texttt{get\_driver}, \texttt{get\_work\_order} and
\texttt{list\_depot\_vehicles} are declared \texttt{read} and
\texttt{assign\_driver} and \texttt{close\_work\_order} are declared
\texttt{write}. The effects come from the registry, not from the verb in the
name --- a distinction that matters, because \texttt{assign} and
\texttt{close} were once absent from a mutating-verb list and thirty writes
were scored as reads. No bulk state export exists, so every argument value
verdict is \textsc{unknown}: the deployment has point lookups and nothing else.

\paragraph{Call A: a legitimate read.} The user asks for the status of vehicle
\texttt{V-0042}; the agent proposes \texttt{get\_vehicle(vehicle\_id="V-0042")}.
Every structural signal is satisfied. The value is grounded --- \texttt{V-0042}
is named in the task text. But \texttt{vehicle\_id} steers where the read
lands, and nothing has confirmed it exists, so autonomy is unavailable:

\begin{center}
{\footnotesize
\begin{tabular}{@{}ll@{}}
\toprule
Step & Outcome \\
\midrule
Argument support & \textsc{unknown} (no coverage declaration) \\
Value grounding & grounded (named in the task) \\
Validation requirement & required (steers the read) \\
Resolver availability & \texttt{get\_vehicle} bound to the role \\
Verdict & \textsc{verify} + \texttt{ResolutionPlan} \\
\midrule
\multicolumn{2}{@{}l@{}}{\emph{validator consulted; entity confirmed present}} \\
\midrule
Re-entry & whole router, fresh observation \\
Final verdict & \textsc{accept} \\
\bottomrule
\end{tabular}}
\end{center}

The plan names one resolver and exactly one argument it may write. The
validator returns a verdict and hands control back; it never performs the
original call and cannot redirect it. Re-entry re-runs the \emph{whole} router,
so every guard applies again to the confirmed call rather than being patched
around. Had no validator been bound to \texttt{vehicle\_id}, the verdict would
be \textsc{abstain}, not \textsc{verify}: promising a verification that cannot
happen is worse than stopping.

\paragraph{Call B: the same confirmation, a different action.} The agent
proposes \texttt{assign\_driver(work\_order\_id="WO-00123",
\texttt{driver\_id}="D-0017")}. Both identifiers confirm against the system of
record exactly as \texttt{V-0042} did. The call is nonetheless held at
\textsc{verify}: the registry declares this tool a write, and the read/write
posture does not trade confirmed state for autonomy on a state change. No
confidence score moves it.

\paragraph{Call C: structurally perfect, wrong tool.} The user asks for a
vehicle status; the agent proposes \texttt{get\_work\_order} with a work-order
identifier that happens to appear elsewhere in the task. Schema valid,
arguments satisfiable, values grounded, nothing tainted --- every signal in
Call A's ladder is satisfied. The sealed BFCL v4 track measured wrong-call
acceptance at 10.9\%. A declared
tool contract does: the task targets a vehicle, the tool acts on a work order,
and the mismatch is established rather than merely unverified. That authority
is implemented and \textbf{unmeasured} (\S\ref{sec:wrongcall-implication}).

% =====================================================================
\section{Negative Results}
\label{sec:negative}

Publication of negative results is essential for scientific
credibility. This section documents active and resolved findings.
Full documentation is in \texttt{NEGATIVE\_RESULTS.md}.

\begin{table}[h]
\centering\small
\caption{Negative results and mitigations.}
\fitcolumn{%
\begin{tabular}{@{}llll@{}}
\toprule
Finding & Sev. & Status & Mitigation \\
\midrule
$\chi$-proxy AUC = 0.39 & Med & Docum. & Exploratory OOD use \\
$T$--$D$ circularity & Med & Resolv. & Structural $T$ \\
Crit.-phase anticorr. & High & Active & Evidence router \\
Full-cov. acc. weak & Med & Active & Selective prediction \\
Oracle diversity part. & Med & Active & Cross-family trio \\
Temp.\ selection falsified & High & \textbf{Active} & Diagnostic only; use calib.\ conf. \\
Evidence signal proxy & Med & Active & NLI benchmark \\
Audit not tamper-proof & Med & By des. & WORM storage \\
\bottomrule
\end{tabular}}
\label{tab:negresults}
\end{table}

\paragraph{Consensus temperature failed fresh-data confirmation (SAP v3).}
On the pre-registered SAP v3 round (1231 fresh BoolQ/TruthfulQA items,
deduplicated against the prior 544-item corpus; frozen Cloudflare
cross-family trio; three-way split dev 493 / risk-cal 370 / test 368, seed
20260727), the consensus-temperature selection signal \emph{failed}
confirmation: test-split AURC 0.0954 vs.\ 0.0664 for a
dev-split-calibrated mean-confidence baseline (paired bootstrap
$\Delta=0.0290$, 95\% CI [0.0119, 0.0503] excluding zero --- temperature
ranks significantly \emph{worse}), and SGR certification at
$r^*=5\%,\delta=0.10$ found no certifiable coverage. The exploratory
temperature advantage seen on the reused 544-item corpus did not transfer;
temperature is retained as a logged diagnostic, not an authoritative
selector. For scale: distinguishing a 90\%-accurate selector from an
80\% full-coverage baseline at conventional thresholds (one-sided
$\alpha=0.05$, power $1-\beta=0.8$) requires on the order of
150--200 \emph{accepted} items---an order of magnitude beyond the
$N_\text{accepted}=18$--25 accepted sets available in the rounds above,
which is why those are reported as directional only. \emph{What
survives:} on the same round, calibrated
confidence-weighted voting scored 87.8\% vs.\ 85.1\% for unweighted
majority on the untouched test split ($N=368$; paired exact McNemar
$p=0.0064$, significant; vs.\ the best single model directional only,
$p=0.077$). Its selection certificates are marginal per arm and no arm
survives family-wise Bonferroni-3, so calibrated confidence is a promising
pre-registered \emph{secondary} that requires its own frozen confirmation
round before any engine integration (SAP v3 \S8 D-3). Full detail:
\path{NEGATIVE_RESULTS.md} \S18 and claim register CLAIM-012 / CLAIM-013.

\paragraph{$\chi$-proxy failure.}
Susceptibility $\chi$ as standalone difficulty predictor: AUC = 0.39
on $N=302$ items---below chance. Exploratory OOD use (not a demonstrated
detector; the live variant is near-constant unless internal clamps bind):
$\chi > 1.45$ (97th percentile) $\Rightarrow$ \textsc{Escalate} with
\texttt{escalate\_adversarial}.

\paragraph{Critical-phase trust anticorrelation.}
Q4 high-trust: 50\% correct; Q1 low-trust: 75\% correct.
Confirmed on $N=32$ real-oracle items and replicated qualitatively
on augmented $N=511$ simulated dataset. No trust threshold separates
correct from incorrect critical-phase items. \emph{Consequence:}
evidence routing---not trust scoring---is the required mechanism for
critical-phase decisions.

\paragraph{$T$--$D$ circularity (resolved).}
Legacy temperature estimator weighted $D$ at 18\% of $T$, creating
a feedback loop ($D \rightarrow T \rightarrow F \rightarrow V$, all
depending on $D$). Resolved by a prompt-only structural temperature estimator, which
computes $T$ without reference to $D$. The estimator itself was
subsequently withdrawn from the method (\S\ref{sec:method}): it failed
its pre-registered fresh-data confirmation, so this circularity fix is
recorded here as history rather than as a live mechanism.

% =====================================================================
\section{Threats to Validity}
\label{sec:threats}

\textbf{Internal validity:} Trust threshold ($\tau^* = 0.2032$,
18th-percentile neg-temperature on 80\% training split) is derived
in-sample; the retired-Groq held-out evaluation ($n=108$, 20\% stratified,
$p=1.45\times10^{-5}$) supported the threshold, but the SAP v3 fresh-data
round did not confirm the temperature signal it selects on
(\S\ref{sec:negative}); 13.8\% author-curated benchmark items (15.1\%
including the adversarial-curated set); synthetic tool-call adversarial
patterns.

\textbf{External validity:} Model coverage is mixed across rounds: the
earlier QA and tool-call rounds used a Groq-hosted three-LLaMA swarm
(one weight family, partial diversity), while the 2026-07 clean round,
the AgentHarm external validation, and the real-LLM baselines use a
Cloudflare Workers AI cross-family trio (Meta LLaMA, Alibaba Qwen,
Mistral). Results on other providers, closed-source models, or quantized
variants may still differ; the benchmark phase distribution may not match
deployment reality; and MultiNLI is a proxy for field evidence retrieval.

\textbf{Construct validity:} The thermodynamic terminology of earlier
versions has been withdrawn (\S\ref{sec:method}). ``Trust'' is a derived scalar, not a frequency probability
of per-item correctness. Utility scores are designed task weights, not
real-world costs.

\textbf{Statistical validity:} 95\% Wilson confidence intervals are
reported for key accuracy figures; Wilson intervals are fixed-$N$
constructions and are not valid under the continuous monitoring that the
REM-020 longitudinal gate performs (optional stopping). A time-uniform
confidence sequence
\cite{ville1939etude,darling1967confidence,howard2021timeuniform,ramdas2023game}
is therefore reported alongside for the gate's cycle-level false-accept
indicator: $0/168$ cycles gives a 95\% time-uniform upper bound of
4.72\% (\path{results/far_confidence_sequence_v1.json}, claim register
CLAIM-011). The sequence is wider than Wilson at any fixed $N$: the
price of validity under data-dependent stopping.

% =====================================================================
\section{Safety, Ethics, and Governance}
\label{sec:safety}

REMORA is designed to \emph{reduce} unsafe autonomous AI actions.
\textsc{Accept} means assurance conditions are met, not that the action
is correct. Human reviewers receiving \textsc{Escalate} retain full
accountability for outcomes. The oracle swarm routes sensitive
operational data through external APIs; data classification requirements
must be evaluated before deployment. The adversarial firewall detects
known injection patterns but is not a complete defense against
sophisticated adversarial inputs.

% =====================================================================
\section{Reproducibility}
\label{sec:repro}

\textbf{Repository:} \url{https://github.com/darklordVirtual/REMORA-research}.
\textbf{Release:} \texttt{\paperversion} (this paper is versioned in lockstep
with the repository release tag).
\textbf{Python:} 3.11+.

\begin{lstlisting}
python -m pip install -e ".[dev]"
make audit        # lint + tests + claim consistency
make benchmark    # all deterministic benchmarks
make credibility-pack
\end{lstlisting}

Live oracle experiments require \texttt{GROQ\_API\_KEY} (earlier rounds)
or \texttt{CLOUDFLARE\_API\_TOKEN} (Cloudflare Workers AI, the 2026-07
clean round). Oracle model versions may change; results are sensitive to
model version updates. Stored artifact experiments are fully deterministic.

\paragraph{License.}
This preprint is distributed as part of the REMORA Licensed Work under
the Business Source License 1.1 (source-available, not open source).
Commercial use requires a separate REMORA Commercial License from the
Licensor, Stian Skogbrott (\texttt{support@luftfiber.no}). Third-party
datasets and materials remain under their own terms; see
\path{THIRD_PARTY_NOTICES.md}.

\textbf{Claim binding:} headline numbers are bound to the machine-readable
claim register (\path{docs/assurance/claim_register_v1.yaml}), the single
authoritative source, through the Markdown source of this paper
(\path{paper/remora_paper.md}): the claim-provenance and claim-consistency
CI gates fail on divergence between that source and the register. This
LaTeX manuscript and the compiled PDF are kept in lockstep with the
Markdown source by \path{scripts/check_paper_sync.py} (version, revision,
and guarded process claims); numeric drift in the typeset copy is caught by
review rather than by an automated gate on the \texttt{.tex}. Each \texttt{review-*}
tag triggers a workflow that regenerates a commit-bound test attestation
and a governance pack whose \path{manifest.json} records the exact commit
SHA and per-file SHA-256 checksums, uploaded as one
\texttt{remora-review-handoff-<sha>} artifact.

% =====================================================================

% TEE content moved to Appendix~\ref{app:tee}.

\section{Conclusion}
\label{sec:conclusion}

We presented REMORA, a policy-gated assurance architecture for governing
agentic AI actions before they execute. Three results carry the paper,
and the third is the one we would most like the field to take up.

\textbf{Autonomous safety needs a deterministic floor.} The
0\% unsafe-acceptance result on the adversarial tool-call benchmark (700 synthetic rows; effective $N=70$ clusters;
cluster-level Wilson CI [0.0\%, 5.2\%]; baselines 1.4\%, difference not
significant at $p=0.50$; the significant gains are utility $+0.456$ and
accuracy) is attributable \emph{entirely} to deterministic hard blocks,
not to the multi-oracle consensus machinery. We state this against our
own interest: the probabilistic architecture that gives the system its
name is not what produces its headline safety number. It contributes
routing quality, and that is a different claim requiring different
evidence.

\textbf{Unknown state is a resolution problem.} Absence without declared
completeness is \textsc{unknown}, not \textsc{unsupported}; a
\textsc{verify} that cannot name a bounded lookup degrades to
\textsc{abstain}; and a resolved value re-enters the whole router rather
than patching the prior decision. Under study conditions this recovers
read utility from 0\% to 100\% in the all-\textsc{unknown} regime with
zero corrupt-identifier acceptances, zero cross-tenant lookups, and
writes held at \textsc{verify} throughout.

\textbf{The routing behaviour transfers.} A sealed BFCL v4 track evaluated
the current pipeline once and met all five pre-registered targets:
wrong-call acceptance \textbf{10.9\%} (28/258), irrelevant-tool refusal
100.0\%, unobtainable-input refusal 99.0\%, obtainable-input verification
97.0\%, and required-UNKNOWN autonomous acceptance 0.0\%. Labelled routing
accuracy was 91.2\%. Structural validity and semantic correctness remain
distinct properties, and the isolated effect of semantic task-contract
matching is not established by this track.

Key negative findings retained in full: trust anticorrelates with
correctness in the critical phase; the $\chi$-proxy fails as a difficulty
predictor (AUC $=0.39$); a pre-registered fresh-data round falsified the
consensus-temperature selection signal, which is retained only as a
diagnostic; evidence retrieval is proxy-based; entropy uses a
token-fingerprint heuristic rather than semantic entropy over NLI
clusters; and the AgentHarm run blocked every benign counterpart as well
as every harmful one.

REMORA is a research-grade prototype, \texttt{SHADOW\_ONLY}, with independent
external replication and production field evidence pending. We offer it as a
reference architecture and empirical baseline. The broader warning remains:
a gate that reasons only about action well-formedness cannot establish that
the action serves the user's goal.

% =====================================================================
\section*{Acknowledgements}
Benchmarks use publicly available datasets (BoolQ, TruthfulQA, MultiNLI,
ARC-Challenge, MMLU-Pro). Oracle inference via the Groq API (earlier
rounds) and Cloudflare Workers AI (the 2026-07 clean round). Open Policy
Agent used for policy evaluation. Negative results documented in full in
\texttt{NEGATIVE\_RESULTS.md}; the authors thank all reviewers who pointed
out scope corrections and limitations. We thank contributor Erlend Barli
(\texttt{github.com/erlendbarli}) for contributions to the project
(see \texttt{CONTRIBUTORS.md}).

% =====================================================================
\section*{Use of AI-Assisted Development Tools}
Generative AI tools were used as assistive tools during the development of
this research prototype. Their use included brainstorming, code drafting
and refactoring, test scaffolding, documentation editing, language
revision, and support in locating and summarising relevant literature.

The author defined the research questions, system architecture, policy
semantics, evaluation criteria, experimental scope, and all reported
claims. All source code, tests, benchmark artifacts, citations, numerical
results, and conclusions were reviewed and verified by the author before
inclusion. Generative AI tools were not listed as authors and were not
relied upon as independent sources of experimental evidence.

Where AI-assisted output informed implementation or prose, the final
responsibility for correctness, reproducibility, and interpretation remains
with the human author. A full disclosure is available in
\texttt{docs/AI\_USE.md}.

% =====================================================================
\bibliographystyle{plainnat}
\begin{thebibliography}{99}
\bibitem{kuhn2023semantic}
Kuhn, L., Gal, Y., \& Farquhar, S.\ (2023).
Semantic uncertainty: Linguistic invariances for uncertainty estimation in
natural language generation.
In \emph{Proceedings of the International Conference on Learning Representations
(ICLR 2023)}.
arXiv:2302.09664.

\bibitem{kirsch2024prover}
Kirchner, J.~H., Chen, Y., Edwards, H., Leike, J., McAleese, N.,
\& Burda, Y.\ (2024).
Prover-verifier games improve legibility of LLM outputs.
\emph{arXiv:2407.13692}.

\bibitem{angelopoulos2022crc}
Angelopoulos, A.\ N., Bates, S., Fisch, A., Lei, L., \& Schuster, T.\ (2022).
Conformal risk control.
\emph{arXiv:2208.02814}.

\bibitem{tibshirani2019conformal}
Tibshirani, R.\ J., Barber, R.\ F., Cand\`es, E.\ J., \& Ramdas, A.\ (2019).
Conformal prediction under covariate shift.
\emph{NeurIPS 2019}.

\bibitem{andriushchenko2024agentharm}
Andriushchenko, M., et al.\ (2024).
AgentHarm: A benchmark for measuring harmfulness of LLM agents.
\emph{arXiv:2410.09024}.

\bibitem{dong2025tee}
Dong, B.\ \& Wang, Q.\ (2025).
Evaluating the performance of the DeepSeek model in confidential computing
environment.
\emph{arXiv:2502.11347}.

\bibitem{wang2023self}
Wang, X., et al.\ (2023).
Self-consistency improves chain of thought reasoning in language models.
\emph{ICLR 2023}.

\bibitem{du2023improving}
Du, Y., Li, S., Torralba, A., Tenenbaum, J., \& Mordatch, I.\ (2024).
Improving factuality and reasoning in language models through multiagent debate.
\emph{ICML 2024}.

\bibitem{zheng2023judging}
Zheng, L., et al.\ (2023).
Judging LLM-as-a-judge with MT-Bench and Chatbot Arena.
\emph{NeurIPS 2023}.

\bibitem{cobbe2021training}
Cobbe, K., et al.\ (2021).
Training verifiers to solve math word problems.
\emph{arXiv:2110.14168}.

\bibitem{chow1970optimum}
Chow, C.\ K.\ (1970).
On optimum recognition error and reject tradeoff.
\emph{IEEE Transactions on Information Theory, 16}(1), 41--46.

\bibitem{elyaniv2010selective}
El-Yaniv, R.\ \& Wiener, Y.\ (2010).
On the foundations of noise-free selective classification.
\emph{Journal of Machine Learning Research, 11}, 1605--1641.

\bibitem{lakshminarayanan2017simple}
Lakshminarayanan, B., Pritzel, A., \& Blundell, C.\ (2017).
Simple and scalable predictive uncertainty estimation using deep ensembles.
\emph{Advances in Neural Information Processing Systems (NeurIPS 2017)}.

\bibitem{kuncheva2003measures}
Kuncheva, L.\ I.\ \& Whitaker, C.\ J.\ (2003).
Measures of diversity in classifier ensembles and their relationship with
the ensemble accuracy.
\emph{Machine Learning, 51}(2), 181--207.

\bibitem{geifman2017selective}
Geifman, Y.\ \& El-Yaniv, R.\ (2017).
Selective classification for deep neural networks.
\emph{NeurIPS 2017}.

\bibitem{kadavath2022language}
Kadavath, S., et al.\ (2022).
Language models (mostly) know what they know.
\emph{arXiv:2207.05221}.

\bibitem{angelopoulos2021gentle}
Angelopoulos, A.\ \& Bates, S.\ (2021).
A gentle introduction to conformal prediction and distribution-free uncertainty quantification.
\emph{arXiv:2107.07511}.

\bibitem{guo2017calibration}
Guo, C., et al.\ (2017).
On calibration of modern neural networks.
\emph{ICML 2017}.

\bibitem{endsley1995toward}
Endsley, M.\ (1995).
Toward a theory of situation awareness in dynamic systems.
\emph{Human Factors, 37}(1), 32--64.

\bibitem{kelly1998goal}
Kelly, T.\ (1998).
Arguing safety: A systematic approach to managing safety cases.
PhD thesis, University of York.

\bibitem{bloomfield2010safety}
Bloomfield, R.\ \& Bishop, P.\ (2010).
Safety and assurance cases: Practice, trends and challenges.
\emph{SAFECOMP 2010}.

\bibitem{lin2022truthfulqa}
Lin, S., et al.\ (2022).
TruthfulQA: Measuring how models mimic human falsehoods.
\emph{ACL 2022}.

\bibitem{clark2019boolq}
Clark, C., et al.\ (2019).
BoolQ: Exploring the surprising difficulty of natural yes/no questions.
\emph{NAACL 2019}.

\bibitem{williams2018broad}
Williams, A., et al.\ (2018).
A broad-coverage challenge corpus for sentence understanding through inference.
\emph{NAACL 2018}.

\bibitem{euaiact2024}
European Parliament \& Council of the European Union.\ (2024).
Regulation (EU) 2024/1689 of the European Parliament and of the Council
(Artificial Intelligence Act).
\emph{Official Journal of the European Union}, L~2024/1689.

\bibitem{opa}
Open Policy Agent.\ (2024).
\url{https://www.openpolicyagent.org/}.

\bibitem{bjoru2026causal}
Bjøru, A.~R.\ (2026).
Causal Post-hoc Explainable AI.
PhD thesis, Norwegian University of Science and Technology (NTNU).
ISBN~978-82-353-0022-5.

\bibitem{chennabasappa2025llamafirewall}
Chennabasappa, S., et al.\ (2025).
LlamaFirewall: An open source guardrail system for building secure AI agents.
\emph{arXiv:2505.03574}.

\bibitem{choi2026membrane}
Choi, M., Yang, S., Kim, D., Kim, S., Son, J., Lee, Y., Choo, J.,
\& Kwak, Y.\ (2026).
Membrane: A self-evolving contrastive safety memory for LLM agent defense.
\emph{arXiv:2606.05743}.

\bibitem{corsi2021formal}
Corsi, D., Marchesini, E., \& Farinelli, A.\ (2021).
Formal verification of neural networks for safety-critical tasks in deep
reinforcement learning.
\emph{UAI 2021}, PMLR 161:333--343.

\bibitem{darling1967confidence}
Darling, D.~A.\ \& Robbins, H.\ (1967).
Confidence sequences for mean, variance, and median.
\emph{Proceedings of the National Academy of Sciences, 58}(1), 66--68.

\bibitem{debenedetti2024agentdojo}
Debenedetti, E., et al.\ (2024).
AgentDojo: A dynamic environment to evaluate prompt injection attacks and
defenses for LLM agents.
\emph{NeurIPS 2024 Datasets \& Benchmarks}.

\bibitem{debenedetti2025camel}
Debenedetti, E., et al.\ (2025).
Defeating prompt injections by design (CaMeL).
\emph{arXiv:2503.18813}.

\bibitem{farquhar2024detecting}
Farquhar, S., Kossen, J., Kuhn, L., \& Gal, Y.\ (2024).
Detecting hallucinations in large language models using semantic entropy.
\emph{Nature, 630}, 625--630.

\bibitem{ge2026governance}
Ge, Y.\ (2026).
Governance architecture for autonomous agent systems: Threats, framework,
and engineering practice.
\emph{arXiv:2603.07191}.

\bibitem{greenblatt2024aicontrol}
Greenblatt, R., Shlegeris, B., Sachan, K., \& Roger, F.\ (2024).
AI control: Improving safety despite intentional subversion.
\emph{ICML 2024}, PMLR 235:16295--16336.

\bibitem{guldimann2024complai}
Guldimann, P., et al.\ (2024).
COMPL-AI framework: A technical interpretation and LLM benchmarking suite
for the EU Artificial Intelligence Act.
\emph{arXiv:2410.07959}.

\bibitem{howard2021timeuniform}
Howard, S.~R., Ramdas, A., McAuliffe, J., \& Sekhon, J.\ (2021).
Time-uniform, nonparametric, nonasymptotic confidence sequences.
\emph{Annals of Statistics, 49}(2), 1055--1080.

\bibitem{iec61511}
International Electrotechnical Commission.\ (2016).
IEC 61511-1:2016, Functional safety: Safety instrumented systems for the
process industry sector.
IEC.

\bibitem{inan2023llamaguard}
Inan, H., et al.\ (2023).
Llama Guard: LLM-based input-output safeguard for human-AI conversations.
\emph{arXiv:2312.06674}.

\bibitem{mohri2024conformal}
Mohri, C.\ \& Hashimoto, T.\ (2024).
Language models with conformal factuality guarantees.
\emph{ICML 2024}.

\bibitem{mu2024rules}
Mu, T., Helyar, A., et al.\ (2024).
Rule based rewards for language model safety.
\emph{NeurIPS 2024}.

\bibitem{norsok2021d010}
Standards Norway.\ (2021).
NORSOK D-010:2021, Well integrity in drilling and well operations.
Standards Norway.

\bibitem{pearl2009causality}
Pearl, J.\ (2009).
\emph{Causality: Models, Reasoning, and Inference} (2nd ed.).
Cambridge University Press.

\bibitem{raji2022outsider}
Raji, I.~D., Xu, P., Honigsberg, C., \& Ho, D.~E.\ (2022).
Outsider oversight: Designing a third party audit ecosystem for AI
governance.
\emph{AIES '22}. arXiv:2206.04737.

\bibitem{ramdas2023game}
Ramdas, A., Gr\"unwald, P., Vovk, V., \& Shafer, G.\ (2023).
Game-theoretic statistics and safe anytime-valid inference.
\emph{Statistical Science, 38}(4), 576--601.

\bibitem{rebedea2023nemo}
Rebedea, T., et al.\ (2023).
NeMo Guardrails: A toolkit for controllable and safe LLM applications with
programmable rails.
\emph{EMNLP 2023 System Demonstrations}.

\bibitem{ruan2024toolemu}
Ruan, Y., et al.\ (2024).
Identifying the risks of LM agents with an LM-emulated sandbox (ToolEmu).
\emph{ICLR 2024}.

\bibitem{shafer2008tutorial}
Shafer, G.\ \& Vovk, V.\ (2008).
A tutorial on conformal prediction.
\emph{Journal of Machine Learning Research, 9}, 371--421.

\bibitem{shamsujjoha2024taxonomy}
Shamsujjoha, M., Lu, Q., Zhao, D., \& Zhu, L.\ (2024).
A taxonomy of multi-layered runtime guardrails for designing foundation
model-based agents: Swiss cheese model for AI safety by design.
\emph{arXiv:2408.02205}.

\bibitem{sharma2025constitutional}
Sharma, M., et al.\ (2025).
Constitutional classifiers: Defending against universal jailbreaks across
thousands of hours of red teaming.
\emph{arXiv:2501.18837}.

\bibitem{shi2025progent}
Shi, T., He, J., Wang, Z., Li, H., Wu, L., Guo, W., \& Song, D.\ (2025).
Progent: Securing AI agents with privilege control.
\emph{arXiv:2504.11703}.

\bibitem{ville1939etude}
Ville, J.\ (1939).
\emph{\'Etude critique de la notion de collectif.}
Gauthier-Villars.

\bibitem{vovk2005algorithmic}
Vovk, V., Gammerman, A., \& Shafer, G.\ (2005).
\emph{Algorithmic Learning in a Random World.}
Springer.

\bibitem{wang2023lora}
Wang, X., Aitchison, L., \& Rudolph, M.\ (2023).
LoRA ensembles for large language model fine-tuning.
\emph{arXiv:2310.00035}.

\bibitem{wang2025agentspec}
Wang, H., Poskitt, C.~M., et al.\ (2025).
AgentSpec: Customizable runtime enforcement for safe and reliable LLM
agents.
\emph{arXiv:2503.18666}.

\bibitem{xiang2025guardagent}
Xiang, Z., et al.\ (2025).
GuardAgent: Safeguard LLM agents by a guard agent via knowledge-enabled
reasoning.
\emph{ICML 2025}.

\bibitem{yadkori2024abstention}
Yadkori, Y.~A., Kuzborskij, I., et al.\ (2024).
Mitigating LLM hallucinations via conformal abstention.
\emph{arXiv:2405.01563}.

\bibitem{yang2026agenttrust}
Yang, C.\ (2026).
AgentTrust: A self-improving trust layer for AI-agent actions.
\emph{arXiv:2606.08539}.

\bibitem{yao2024taubench}
Yao, S., Shinn, N., Razavi, P., \& Narasimhan, K.\ (2024).
$\tau$-bench: A benchmark for tool-agent-user interaction in real-world
domains.
\emph{arXiv:2406.12045}.

\bibitem{yuan2024rjudge}
Yuan, T., et al.\ (2024).
R-Judge: Benchmarking safety risk awareness for LLM agents.
\emph{Findings of EMNLP 2024}.

\bibitem{zhan2024injecagent}
Zhan, Q., Liang, Z., et al.\ (2024).
InjecAgent: Benchmarking indirect prompt injections in tool-integrated LLM
agents.
\emph{Findings of ACL 2024}.

\bibitem{zhang2024fidelity}
Zhang, M., Huang, M., Shi, R., Guo, L., Peng, C., Yan, P., Zhou, Y.,
\& Qiu, X.\ (2024).
Calibrating the confidence of large language models by eliciting fidelity.
\emph{arXiv:2404.02655}.

\end{thebibliography}

% =====================================================================
\appendix

\section{Mathematical Notation Table}
\label{app:math}

\begin{table}[h]
\centering\small
\caption{Mathematical notation.}
\fitcolumn{%
\begin{tabular}{@{}lll@{}}
\toprule
Symbol & Definition & Domain \\
\midrule
$\Oracles$ & Oracle set & $n \geq 2$ endpoints (design intent $\geq 3$ distinct families) \\
$\OraclesV$ & Valid oracle set & $\OraclesV \subseteq \Oracles$ \\
$\phi(o)$ & Canonical verdict & (polarity, hash, mag., tags) \\
$\rho(a,b)$ & Pairwise agreement rate & $[0,1]$, window 200 \\
$w(o)$ & Diversity weight & normalized \\
$\hat{p}(v)$ & Weighted support & $[0,1]$, sums to 1 \\
$H$ & Shannon entropy & bits $\geq 0$ \\
$D$ & Dissensus & $[0,1]$ \\
$k$ & Active verdict count & $\geq 1$ \\
$\eta$ & Consensus concentration & $[0,1]$ \\
$\chi$ & Susceptibility & $[0, \infty)$ \\
$\lambda$ & Coupling constant & 0.3 (paper config; library default 1.0) \\
$\trust$ & Trust score & $[0,1]$ \\
$h_{\rm bound}$ & Heuristic false-consensus risk proxy & $[0,1]$ \\
$w_{\rm phase}$ & Phase weight & $\{1.0, 0.5, 0.1\}$ \\
$g$ & Gate outcome & ACCEPT/VERIFY/ABS./ESC. \\
$\Gate$ & Gate function & $(q,a,c) \to g$ \\
\bottomrule
\end{tabular}}
\end{table}

\section{DecisionEnvelope v2 Schema Summary}
\label{app:schema}

Top-level blocks: \texttt{request}, \texttt{assessment},
\texttt{gate}, \texttt{reviewer\_context}, \texttt{follow\_up},
\texttt{history}, \texttt{policy\_learning}, \texttt{audit}.
Full JSON schema is in the Markdown paper and repository
(\texttt{remora/policy/report.py}).

\section{Negative Results Archive}
\label{app:negres}

The complete, versioned record of negative and resolved findings ---
including the resolved items R1 (iteration damage), R2 (Stage-3b critique
impact) and R7 ($T$--$D$ circularity) --- lives in \path{NEGATIVE_RESULTS.md}
in the repository. It is not duplicated here: the main-body Negative Results
section carries the findings that shape the architecture; this appendix is a
single pointer to the full archive so the two do not drift.

\section{Implementation Status Table}
\label{app:impl}

Status vocabulary: \emph{Integrated} = invoked by the primary
\texttt{/v1/assess} path; \emph{Offline} = implemented and evaluated with
committed artifacts but not invoked by the primary API; \emph{Adapter} =
implemented and parity-tested, available as an integration option;
\emph{Planned} = not implemented.

\begin{table}[h]
\centering\small
\caption{Component status: implemented vs.\ integrated in \texttt{/v1/assess}.}
\fitcolumn{%
\begin{tabular}{@{}llll@{}}
\toprule
Component & Impl. & Integrated & Notes \\
\midrule
Oracle fan-out & Yes & Yes & ThreadPoolExecutor \\
Canonicalization ($\phi$) & Yes & Yes & NFKC, Semantic Entropy \cite{kuhn2023semantic} \\
Correlation weighting & Yes & Per-request & Fresh engine per API request \\
$H$, $D$, $\eta$, $\trust$ & Yes & Yes & \texttt{thermodynamics.py} \\
$T$, $F$ (withdrawn, \S\ref{sec:negative}) & Yes & Yes & \texttt{thermodynamics.py} \\
$V(t)$ session stability (withdrawn, \S\ref{sec:negative}) & Yes & Yes & Abort on $\Delta V$ \\
Phase classification & Yes & Yes & 3-class \\
Policy engine (7 headline controls: 3 unconditional + 4 conditional) & Yes & Yes & Full gate inventory in code \\
OPA/Rego adapter & Adapter & No & API calls Python engine directly \\
Governance Intelligence & Offline & No & Optional pre-policy helper \\
\texttt{PhaseAwareGuardrail} & Offline & No & Committed $N=544$ artifact \\
Covariate-shift CRC & Offline & No & Nominal bound (heuristic weights) \\
PVD-inspired agreement score & Offline & No & Post-hoc heuristic \\
Evidence router (oracle proxy) & Yes & Yes & Proxy signal \\
Evidence router (semantic) & Planned & No & Interface ready \\
DecisionEnvelope v2 & Yes & Yes & Attribute-frozen, not deep-frozen \\
SHA-256 hash-chain & Yes & API layer & HMAC only with signing key \\
Execution grant path & Yes & Separate router & Dispatches via \texttt{GovernedToolDispatcher} (lease-bound; CAP-013) \\
Human review workflow & Yes & Partial & Demo UI only \\
WORM audit storage & Planned & No & External dependency \\
ZK assurance proofs & Planned & No & Stubbed \\
\bottomrule
\end{tabular}}
\end{table}

\section{Conformal Risk Control Under Covariate Shift (offline component)}
\label{app:crc}


Standard Mondrian conformal calibration assumes that calibration and
test samples are \emph{exchangeable} - drawn i.i.d.\ from the same
distribution.  This assumption is violated in the critical phase: the
joint distribution $(\tau, \text{correct})$ differs structurally
between ordered-phase calibration items and critical-phase test items
due to the trust-score inversion documented above.  Empirically,
applying standard conformal to critical-phase items yields 100\%
observed risk and zero coverage.

REMORA addresses this with a \textbf{weighted empirical selective
router} \emph{inspired by} Conformal Risk Control (CRC,
\cite{angelopoulos2022crc}) but --- important --- \emph{not
implementing the CRC procedure} (see below).  Each calibration item $i$
receives an importance weight $w_i = p_{\text{test}}(x_i) /
p_{\text{cal}}(x_i)$.  For REMORA's phase shift, we use the
hand-tuned estimate

\begin{equation}
  w_i = \begin{cases} 1.0 & \text{if phase}(i) = p_{\text{test}} \\
                      \beta & \text{otherwise} \end{cases}
  \qquad (\beta = 0.10)
  \label{eq:phase_weights}
\end{equation}

The selector picks the \emph{smallest} threshold $\hat{\lambda}$ (i.e.\
maximum coverage) such that the weighted empirical risk
$\bar{L}(\lambda) = \sum_i \tilde{w}_i \ell_i(\lambda) \leq \alpha$,
with $\tilde{w}_i = w_i / \sum_j w_j$ normalised over the calibration
items.

\begin{theorem}[CRC Guarantee, Angelopoulos et al.\ 2022 --- stated for
context; not instantiated by the implementation]
For any monotone loss $\ell : [0,1] \to [0,B]$, target risk
$\alpha \in (0,B)$, and correctly-specified importance weights:
\[
  \mathbb{E}[L(\hat{\lambda})] \leq \alpha + \frac{B}{n+1}
\]
where $n$ is the calibration set size.
\end{theorem}

\textbf{Non-applicability.} The implemented selector does not
instantiate this theorem, for two structural reasons no weight
specification can repair (external review R3-01): (a)~its threshold
rule applies the raw constraint $\bar{L}(\lambda) \leq \alpha$ with no
$B/(n+1)$ finite-sample term and no test-point weight $w_{n+1}$;
(b)~the controlled quantity --- the conditional error rate among
accepted items --- is not monotone in $\lambda$, so Theorem~1 does not
attach even with valid density-ratio weights. In addition,
$\beta = 0.10$ is a hand-tuned constant, not a density ratio
\cite{tibshirani2019conformal}. All quality evidence for this component is
empirical (repeated-split holdout results).

The component is implemented as
\path{remora.selective.crc.WeightedEmpiricalSelectiveRouter}
(formerly \texttt{CovariateShiftCRC}, retained as an alias) with a
\texttt{fit(scores, labels, phases, target\_phase)} interface
compatible with \texttt{ConformalPhaseGuardrail}.  A \texttt{CRCReport}
dataclass exposes \texttt{finite\_sample\_slack} ($= 1/(n_{\text{cal}}+1)$)
and \texttt{nominal\_risk\_bound} ($= \alpha + \text{slack}$) as
informational reference values (not applicable bounds), plus
\texttt{holdout\_risk} (plain unweighted holdout count). The module is
fully tested (45 unit tests) and is an
offline research component, not invoked by the reference API.



\section{Optional TEE-Based Audit Hardening (Tamper-Resistant)}
\label{app:tee}

\emph{This appendix describes tamper-resistant (not tamper-proof) audit hardening
via hardware attestation \cite{dong2025tee}. Implementation status: protocol specified; hardware
integration pending. The term ``tamper-resistant'' is used throughout; the
hash-chain in the main system is tamper-detecting only.}

\subsection{Motivation}

The SHA-256 hash-chain in \path{remora/audit/hash_chain.py} detects
post-hoc modification of audit records but does not prevent an adversary with
storage write access from replacing the entire chain. For deployments where
stronger guarantees are required, hardware attestation via a Trusted Execution
Environment (TEE) provides a complementary layer: the \texttt{DecisionEnvelope}
is signed inside a hardware-isolated enclave, and the attestation quote
includes a measurement of the enclave binary, making silent substitution
detectable even if storage is compromised.

This is an optional deployment hardening; it is not required for the core
REMORA governance protocol described in the main paper.

\subsection{Supported TEE Platforms}

Three platforms are in scope for the reference implementation:
\begin{itemize}
  \item \textbf{Intel TDX / SGX} - remote attestation via DCAP quote generation
  \item \textbf{AMD SEV-SNP} - firmware measurement in attestation report
  \item \textbf{ARM TrustZone} - OP-TEE trusted application (mobile / edge)
\end{itemize}

\subsection{DecisionEnvelope Attestation Format}

When TEE hardening is active, each \texttt{DecisionEnvelope} includes an
\texttt{attestation} sub-object:

\begin{lstlisting}
{
  "envelope_id": "env_7f3a...",
  "gate": { "outcome": "ESCALATE", ... },
  "attestation": {
    "tee_platform": "intel_tdx",
    "quote":        "<base64 DCAP quote>",
    "mrenclave":    "<SHA-256 of enclave measurement>",
    "signature":    "<Ed25519 over envelope_id
                     +outcome+timestamp>",
    "timestamp":    "2026-01-15T14:23:07Z"
  }
}
\end{lstlisting}

The \texttt{mrenclave} value pins the exact binary that produced the decision,
allowing auditors to verify that no unauthorised code modification occurred
between decisions.

\subsection{Verification Flow}

\begin{enumerate}
  \item Auditor retrieves envelope and attestation quote from audit log.
  \item Quote is verified against Intel/AMD/ARM attestation service.
  \item \texttt{mrenclave} is compared against the published release measurement.
  \item Ed25519 signature over \texttt{(envelope\_id, outcome, timestamp)} is
        verified with the enclave's ephemeral key, which is bound to the quote.
\end{enumerate}

If any step fails, the audit record is flagged as potentially tampered.

\subsection{Implementation Status}

Protocol specified and data contract defined. Hardware integration (enclave
compilation, DCAP provisioning, OP-TEE TA signing) is a deployment
dependency and is not part of the core source-available release. Stub interfaces
are present in \texttt{remora/audit/} for downstream integration.

\end{document}
